\documentclass[12pt,a4paper]{article}
\usepackage[utf8]{inputenc}
\usepackage[T1]{fontenc}
\usepackage[french]{babel}
\usepackage{newtxtext,newtxmath}
\usepackage{geometry}
\usepackage{amsmath}
\usepackage{amsfonts}
\usepackage{graphicx}
\usepackage{svg}
\usepackage{ifthen}
\usepackage{xstring}
\usepackage{calc}
\usepackage{pgfopts}
\usepackage{tikz}
\usepackage{titlesec}
\usepackage{enumitem}
\usepackage{fancyhdr}
\usepackage{xcolor}
\usepackage{colortbl}
\usepackage{array}
\usepackage{booktabs}
\usepackage{longtable}
\usepackage{parskip}
\usepackage{tcolorbox}
\tcbuselibrary{skins, breakable}
\usepackage{float}
\usepackage{underscore}
\usepackage{multirow}
\usepackage{minted}

\usetikzlibrary{shapes,arrows,positioning,fit,backgrounds}

% --- Layout ---
\geometry{left=2.5cm,right=2.5cm,top=2.5cm,bottom=2.5cm}
\setlength{\headheight}{15pt}

% --- Colors ---
\definecolor{primary}{RGB}{0,51,102}
\definecolor{secondary}{RGB}{200,220,240}
\definecolor{accent}{RGB}{0,102,204}

% --- Typography ---
\titleformat{\section}{\Large\bfseries\color{primary}}{\thesection}{1em}{}
\titleformat{\subsection}{\large\bfseries\color{primary}}{\thesubsection}{1em}{}
\titleformat{\subsubsection}{\normalsize\bfseries\color{primary}}{\thesubsubsection}{1em}{}

% --- Header/Footer ---
\pagestyle{fancy}
\fancyhf{}
\fancyhead[L]{\nouppercase{\leftmark}}
\fancyhead[R]{Rapport de Projet de Fin d'Études}
\fancyfoot[C]{\thepage}
\renewcommand{\headrulewidth}{0.4pt}
\renewcommand{\arraystretch}{1.2}

% --- Diagrams ---
\newcommand{\pumlsvg}[2][0.8\textwidth]{
    \includesvg[width=#1, inkscapelatex=false]{../../diagrams/svg/#2.svg}
}

% --- Project Boxes (Modern Flat Style) ---
\newtcolorbox{projectbox}[1]{
    enhanced,
    breakable,
    colback=primary!5,
    colframe=primary,
    boxrule=0pt,
    leftrule=4pt,
    sharp corners,
    fonttitle=\bfseries,
    coltitle=primary,
    attach title to upper,
    after title={ :\enskip},
    title=#1,
    top=10pt,
    bottom=10pt,
    left=10pt,
    right=10pt
}

% --- Hyperlinks (MUST BE LOADED BEFORE CLEVEREF) ---
\usepackage[
    colorlinks = true,
    linkcolor  = primary,
    urlcolor   = accent,
    citecolor  = primary,
    bookmarks  = true
]{hyperref}

% --- Cross-referencing (MUST BE LOADED AFTER HYPERREF) ---
\usepackage[capitalise, french, noabbrev]{cleveref}

\crefname{figure}{figure}{figures}
\Crefname{figure}{Figure}{Figures}
\crefname{table}{tableau}{tableaux}
\Crefname{table}{Tableau}{Tableaux}
\crefname{section}{section}{sections}
\Crefname{section}{Section}{Sections}


\hypersetup{
    pdftitle   = {Rapport PFE -- Système de gestion des soutenances et des jurys},
    pdfauthor  = {Zaitouni Nourelislam \& Taoudi Ebourki Abdelaziz},
}

\begin{document}

% ── Title Page ─────────────────────────────────────────────────────────────────
\begin{titlepage}
    \centering
    \vspace*{0.5cm}
    \includegraphics[width=0.4\textwidth]{../../assets/logo.png}\\[1cm]
    {\Large \textbf{Faculté des Sciences Ben M'Sick (FSBM)}}\\[0.5cm]
    {\large Université Hassan II de Casablanca}\\[1.5cm]

    \rule{\linewidth}{0.5mm}\\[0.4cm]
    {\huge \bfseries \textcolor{primary}{Rapport de Projet de Fin d'Études}}\\[0.2cm]
    {\Large \bfseries Système de gestion des soutenances et des jurys}\\
    \rule{\linewidth}{0.5mm}\\[1.5cm]

    \begin{minipage}{0.5\textwidth}
        \begin{flushleft}\large
            \emph{Réalisé par :}\\
            \textbf{Zaitouni Nourelislam}\\
            \textbf{Taoudi Ebourki Abdelaziz}
        \end{flushleft}
    \end{minipage}%
    \begin{minipage}{0.45\textwidth}
        \begin{flushright}\large
            \emph{Encadré par :}\\
            \textbf{Pr. Mohammed AIT DAOUD}\\[0.5cm]
            \emph{Co-encadré par :}\\
            \textbf{Dr. Salma SAMMAH}
        \end{flushright}
    \end{minipage}

    \vfill
    \begin{center}
        \textbf{Licence Fondamentale -- MIP}\\[0.3cm]
        Année Universitaire 2025--2026
    \end{center}
\end{titlepage}

% ── Dédicace ───────────────────────────────────────────────────────────────────
\newpage
\thispagestyle{empty}
\vspace*{\fill}
\begin{center}
    \textit{À nos parents, pour leur soutien indéfectible.}\\[0.5cm]
    \textit{À nos professeurs, pour leur encadrement et leur savoir.}\\[0.5cm]
    \textit{À Dr. Salma SAMMAH, pour son accompagnement exceptionnel et ses conseils précieux tout au long de ce projet.}\\[0.5cm]
    \textit{À tous ceux qui ont contribué, de près ou de loin, à la réalisation de ce projet.}
\end{center}
\vspace*{\fill}

% ── Remerciements ──────────────────────────────────────────────────────────────
\newpage
\thispagestyle{empty}
\section*{Remerciements}
\addcontentsline{toc}{section}{Remerciements}

Nos premiers remerciements s'adressent à \textbf{Dr. Salma SAMMAH}, pour son accompagnement technique constant, sa disponibilité exemplaire et la qualité de ses orientations qui ont été déterminantes pour la réalisation de ce projet. Son expertise, ses conseils avisés et ses retours constructifs nous ont permis d'avancer efficacement tout au long du développement.

Nous remercions également \textbf{Pr. Mohammed AIT DAOUD} pour la proposition de ce sujet de projet et son encadrement.

Nos remerciements s'adressent aussi à l'ensemble du corps professoral du département MIP de la Faculté des Sciences Ben M'Sick pour la qualité de la formation dispensée durant ces trois années de licence.

Enfin, nous remercions nos familles et nos proches pour leur patience et leurs encouragements constants.

% ── Résumé / Abstract ──────────────────────────────────────────────────────────
\newpage
\thispagestyle{empty}
\section*{Résumé}
\addcontentsline{toc}{section}{Résumé}

L'organisation des soutenances en milieu universitaire constitue un processus complexe mobilisant de multiples acteurs -- étudiants, encadrants, membres de jury et personnel administratif -- autour de contraintes temporelles, spatiales et réglementaires. Le présent projet vise à concevoir et développer une application web complète pour la gestion centralisée, automatisée et fiable des soutenances et des jurys au sein de la Faculté des Sciences Ben M'Sick. Le système repose sur une architecture moderne à trois couches : un frontend React, un backend Spring Boot et une base de données MySQL. Il couvre l'ensemble du cycle de vie d'une soutenance, de la planification des créneaux et de l'affectation des jurys jusqu'à la saisie des évaluations et la génération automatique des documents administratifs (convocations, procès-verbaux, fiches d'évaluation). Une attention particulière est portée à la détection des conflits horaires et à l'équilibrage de la charge des enseignants. Ce rapport présente l'étude de l'existant, l'analyse des besoins, la conception UML, l'architecture technique et la modélisation de la base de données du système.

\vspace{0.5cm}
\noindent\textbf{Mots-clés :} Gestion de soutenances, planification, jury, évaluation, Spring Boot, React, JWT, UML, système d'information universitaire.

\vspace{1cm}
\section*{Abstract}
\addcontentsline{toc}{section}{Abstract}

Organizing thesis defenses in a university environment is a complex process involving multiple actors -- students, supervisors, jury members, and administrative staff -- around temporal, spatial and regulatory constraints. This project aims to design and develop a complete web application for the centralized, automated and reliable management of defenses and juries at the Faculty of Sciences Ben M'Sick. The system is based on a modern three-tier architecture: a React frontend, a Spring Boot backend, and a MySQL database. It covers the entire lifecycle of a defense, from slot planning and jury assignment to evaluation entry and automatic generation of administrative documents (summons, minutes, evaluation sheets). Special attention is given to conflict detection and teacher workload balancing. This report presents the study of the current system, requirements analysis, UML design, technical architecture, and database modeling.

\vspace{0.5cm}
\noindent\textbf{Keywords:} Defense management, planning, jury, evaluation, Spring Boot, React, JWT, UML, university information system.

% ── Tables ─────────────────────────────────────────────────────────────────────
\newpage
\tableofcontents
\newpage
\listoffigures
\addcontentsline{toc}{section}{Liste des figures}
\listoftables
\addcontentsline{toc}{section}{Liste des tableaux}

% ── Liste des acronymes ──────────────────────────────────────────────────────────
\newpage
\section*{Liste des acronymes}
\addcontentsline{toc}{section}{Liste des acronymes}

\begin{table}[H]
    \centering
    \small
    \begin{tabular}{llp{8cm}}
        \toprule
        \rowcolor{primary!10}
        \textbf{Acronyme} & & \textbf{Définition} \\
        \midrule
        \textbf{ADMIN} & & Administrateur \\
        \addlinespace
        \textbf{API} & & Application Programming Interface -- Interface de programmation applicative \\
        \addlinespace
        \textbf{BDD} & & Base de Données \\
        \addlinespace
        \textbf{CDC} & & Cahier des Charges \\
        \addlinespace
        \textbf{CI/CD} & & Continuous Integration / Continuous Deployment -- Intégration continue / Déploiement continu \\
        \addlinespace
        \textbf{CRUD} & & Create, Read, Update, Delete -- Opérations de base sur les données \\
        \addlinespace
        \textbf{CSV} & & Comma-Separated Values -- Valeurs séparées par virgule (format de fichier) \\
        \addlinespace
        \textbf{CSS} & & Cascading Style Sheets -- Feuilles de style en cascade \\
        \addlinespace
        \textbf{E2E} & & End-to-End -- De bout en bout (tests) \\
        \addlinespace
        \textbf{FSBM} & & Faculté des Sciences Ben M'Sick \\
        \addlinespace
        \textbf{HTML} & & HyperText Markup Language -- Langage de balisage hypertexte \\
        \addlinespace
        \textbf{HTTP} & & HyperText Transfer Protocol -- Protocole de transfert hypertexte \\
        \addlinespace
        \textbf{HTTPS} & & HyperText Transfer Protocol Secure -- HTTP sécurisé \\
        \addlinespace
        \textbf{IDE} & & Integrated Development Environment -- Environnement de développement intégré \\
        \addlinespace
        \textbf{JPA} & & Java Persistence API -- API de persistance Java \\
        \addlinespace
        \textbf{JWT} & & JSON Web Token -- Jeton web JSON (authentification) \\
        \addlinespace
        \textbf{JS} & & JavaScript -- Langage de programmation web \\
        \addlinespace
        \textbf{JSON} & & JavaScript Object Notation -- Format d'échange de données \\
        \addlinespace
        \textbf{MIP} & & Mathématiques, Informatique et Physique -- Filière de licence \\
        \addlinespace
        \textbf{MVC} & & Model-View-Controller -- Modèle-Vue-Contrôleur (architecture logicielle) \\
        \addlinespace
        \textbf{MySQL} & & Système de gestion de base de données relationnelle \\
        \addlinespace
        \textbf{PDF} & & Portable Document Format -- Format de document portable \\
        \addlinespace
        \textbf{PFE} & & Projet de Fin d'Études \\
        \addlinespace
        \textbf{REST} & & Representational State Transfer -- Style d'architecture pour APIs web \\
        \addlinespace
        \textbf{SQL} & & Structured Query Language -- Langage de requête structuré \\
        \addlinespace
        \textbf{TS} & & TypeScript -- Sur-ensemble typé de JavaScript \\
        \addlinespace
        \textbf{TSX} & & TypeScript XML -- TypeScript avec syntaxe JSX pour React \\
        \addlinespace
        \textbf{UML} & & Unified Modeling Language -- Langage de modélisation unifié \\
        \addlinespace
        \textbf{UI} & & User Interface -- Interface utilisateur \\
        \addlinespace
        \textbf{UX} & & User Experience -- Expérience utilisateur \\
        \addlinespace
        \textbf{XML} & & eXtensible Markup Language -- Langage de balisage extensible \\
        \bottomrule
    \end{tabular}
\end{table}

% ═══════════════════════════════════════════════════════════════════════════════
%  INTRODUCTION GÉNÉRALE
% ═══════════════════════════════════════════════════════════════════════════════
\newpage
\section*{Introduction Générale}
\addcontentsline{toc}{section}{Introduction Générale}
\label{sec:introduction_generale}

\subsection*{Contexte général}
\label{intro:contexte}
\addcontentsline{toc}{subsection}{Contexte général}

L'organisation des soutenances en milieu universitaire est une activité critique qui se heurte aujourd'hui à plusieurs défis. La multiplicité des acteurs -- étudiants, encadrants, membres de jury et services administratifs -- nécessite une coordination fine et rigoureuse. Pourtant, les processus traditionnels reposent encore largement sur des tableurs Excel, des échanges d'e-mails massifs et des documents papier, engendrant une perte d'efficacité significative et des erreurs de planification fréquentes.

Dans ce contexte, la faculté des Sciences Ben M'Sick (FSBM) de l'Université Hassan II de Casablanca, à travers son département de Mathématiques, Informatique et Physique (MIP), cherche à moderniser ses outils de gestion académique. L'objectif est de centraliser les données et d'automatiser les flux de travail via une plateforme web dédiée, capable de répondre aux exigences croissantes de qualité et de réactivité.

\subsection*{Problématique}
\label{intro:problematique}
\addcontentsline{toc}{subsection}{Problématique}

La coordination des soutenances dans un département universitaire fait face à une complexité croissante. Les contraintes à respecter sont multiples et souvent interdépendantes. La disponibilité des enseignants, qui interviennent à la fois comme encadrants et comme membres de jury, doit être prise en compte pour éviter les doubles affectations horaires. La réservation des salles, dont les capacités varient, nécessite une allocation cohérente avec le nombre de participants attendus.

Par ailleurs, la production des documents administratifs accompagnant chaque soutenance -- convocations, fiches d'évaluation -- représente une charge de travail récurrente et chronophage. La saisie et le suivi des évaluations, la transmission des notes aux services compétents et la génération des rapports de synthèse complètent ce tableau.

Le cœur de la problématique réside donc dans l'absence d'un outil unifié capable de gérer l'ensemble de ces aspects de manière intégrée et automatisée. L'organisation manuelle actuelle, basée sur des outils génériques non adaptés à cette finalité spécifique, génère des inefficacités, des erreurs et une traçabilité insuffisante.

\subsection*{Objectifs du projet}
\label{intro:objectifs}
\addcontentsline{toc}{subsection}{Objectifs du projet}

\begin{projectbox}{Objectif principal}
    Développer une application web complète et fonctionnelle permettant la gestion centralisée, automatisée et fiable de l'organisation des soutenances et des jurys au sein d'un établissement universitaire.
\end{projectbox}

Plusieurs objectifs spécifiques déclinent cette vision générale :
\begin{itemize}[label=\textcolor{primary}{\textbullet}]
    \item \textbf{Gestion des acteurs académiques :} étudiants, encadrants, membres de jury et responsables pédagogiques.
    \item \textbf{Planification :} gestion des dates, horaires et salles, avec affectation des jurys et détection automatique des conflits.
    \item \textbf{Production documentaire :} génération automatique des convocations, fiches d'évaluation et procès-verbaux.
    \item \textbf{Évaluation :} saisie des notes et observations, enregistrement des délibérations et production d'états récapitulatifs.
\end{itemize}

\subsection*{Méthodologie suivie}
\label{intro:methodologie}
\addcontentsline{toc}{subsection}{Méthodologie suivie}

La réalisation de ce projet a suivi une démarche structurée en plusieurs phases :
\begin{enumerate}[label=\textcolor{primary}{\arabic*}.]
    \item \textbf{Analyse de l'existant :} étude des processus actuels d'organisation des soutenances, identification des lacunes et définition des besoins.
    \item \textbf{Conception :} modélisation UML (cas d'utilisation, diagramme de classes, diagrammes de séquence et d'activité), conception de l'architecture technique et modélisation de la base de données.
    \item \textbf{Réalisation :} implémentation du backend Spring Boot et du frontend React, développement des modules fonctionnels.
    \item \textbf{Tests et validation :} tests fonctionnels, tests d'intégration et validation utilisateur.
\end{enumerate}

\subsection*{Structure du rapport}
\label{intro:structure}
\addcontentsline{toc}{subsection}{Structure du rapport}

Le présent rapport s'articule autour de cinq chapitres :
\begin{itemize}[label=\textcolor{primary}{\textbullet}]
    \item \textbf{Chapitre 1} présente l'organisme d'accueil, la Faculté des Sciences Ben M'Sick, son organisation et le contexte du stage.
    \item \textbf{Chapitre 2} dresse un état des lieux du système actuel, analyse les solutions existantes et justifie les choix technologiques retenus.
    \item \textbf{Chapitre 3} détaille l'analyse des besoins, la modélisation UML et la conception complète du système.
    \item \textbf{Chapitre 4} décrit la réalisation, l'implémentation technique et les difficultés rencontrées.
    \item \textbf{Chapitre 5} présente la stratégie de test, les scénarios de validation et les limites du projet.
\end{itemize}
Le rapport se termine par une conclusion générale dressant le bilan du travail accompli et les perspectives d'évolution.

% ═══════════════════════════════════════════════════════════════════════════════
%  CHAPITRE 1 : PRÉSENTATION DE L'ORGANISME D'ACCUEIL
% ═══════════════════════════════════════════════════════════════════════════════
\newpage
\section{Présentation de l'organisme d'accueil}
\label{ch1}

\subsection{Introduction}
\label{ch1:intro}

Ce chapitre présente l'environnement dans lequel s'est déroulé notre projet de fin d'études. Il décrit la Faculté des Sciences Ben M'Sick (FSBM) -- son histoire, ses missions, sa structure organisationnelle -- ainsi que le département d'accueil et le cadre spécifique du stage. Cette contextualisation permet de mieux comprendre les enjeux institutionnels et pédagogiques qui ont motivé la réalisation de notre système de gestion des soutenances.

\subsection{Présentation générale de l'organisme}
\label{ch1:presentation}

\subsubsection{Historique et fiche signalétique}
\label{ch1:historique}

La Faculté des Sciences Ben M'Sick (FSBM) est l'une des composantes de l'Université Hassan II de Casablanca, l'une des plus grandes universités du Maroc. Située dans l'arrondissement de Ben M'Sick, à Casablanca, elle a été créée pour répondre à la demande croissante de formation scientifique dans la région du Grand Casablanca.

\begin{table}[H]
    \centering
    \small
    \begin{tabular}{lp{10cm}}
        \toprule
        \rowcolor{primary!10}
        \textbf{Information} & \textbf{Détail} \\
        \midrule
        Nom officiel         & Faculté des Sciences Ben M'Sick (FSBM) \\
        Université de rattachement & Université Hassan II de Casablanca \\
        Adresse              & Avenue Commandant Driss El Harti, Ben M'Sick, Casablanca \\
        Année de création    & 1997 \\
        Type                 & Établissement public d'enseignement supérieur \\
        Domaines             & Sciences exactes, Sciences et technologies \\
        \bottomrule
    \end{tabular}
    \caption{Fiche signalétique de la FSBM}
    \label{tab:fsbm_fiche}
\end{table}

\subsubsection{Missions et activités}
\label{ch1:missions}

La FSBM a pour missions principales :
\begin{itemize}[label=\textcolor{primary}{\textbullet}]
    \item \textbf{La formation initiale :} dispenser un enseignement supérieur de qualité dans les disciplines scientifiques fondamentales et appliquées aux niveaux Licence, Master et Doctorat.
    \item \textbf{La recherche scientifique :} mener des activités de recherche au sein de ses laboratoires et équipes de recherche dans divers domaines des sciences exactes.
    \item \textbf{L'insertion professionnelle :} préparer les étudiants à leur insertion dans le monde professionnel à travers des formations adaptées aux besoins du marché du travail.
    \item \textbf{Le rayonnement académique :} contribuer au développement scientifique et technologique de la région et du pays.
\end{itemize}

\subsubsection{Secteurs d'intervention}
\label{ch1:secteurs}

La faculté couvre un large spectre de disciplines scientifiques :
\begin{itemize}[label=\textcolor{primary}{\textbullet}]
    \item \textbf{Mathématiques :} algèbre, analyse, probabilités, statistiques, mathématiques appliquées.
    \item \textbf{Informatique :} développement logiciel, bases de données, réseaux, intelligence artificielle.
    \item \textbf{Physique :} physique fondamentale, électronique, mécanique, énergétique.
    \item \textbf{Chimie :} chimie organique, inorganique, chimie analytique, biochimie.
    \item \textbf{Sciences de la vie :} biologie moléculaire, biotechnologie, écologie.
    \item \textbf{Sciences de la Terre :} géologie, géophysique, géomatique.
\end{itemize}

\subsection{Structure organisationnelle}
\label{ch1:organisation}

\subsubsection{Organigramme}
\label{ch1:organigramme}

La FSBM est administrée par un Doyen, assisté par des Vice-Doyens et un Secrétaire Général. La structure s'articule autour de plusieurs organes et entités :

\begin{itemize}[label=\textcolor{primary}{\textbullet}]
    \item \textbf{Conseil de faculté :} instance délibérante composée d'enseignants-chercheurs, d'étudiants et de représentants administratifs.
    \item \textbf{Commissions pédagogiques :} chargées de l'organisation et du suivi des filières de formation.
    \item \textbf{Départements pédagogiques :} entités regroupant les enseignants-chercheurs par discipline.
    \item \textbf{Laboratoires de recherche :} structures dédiées à la recherche scientifique.
    \item \textbf{Services administratifs :} scolarité, ressources humaines, finances, bibliothèque.
\end{itemize}

\subsubsection{Description des départements}
\label{ch1:departements}

La faculté est structurée en plusieurs départements pédagogiques, chacun responsable d'une ou plusieurs filières de formation :

\begin{itemize}[label=\textcolor{primary}{\textbullet}]
    \item \textbf{Département de Mathématiques :} propose les licences en Mathématiques Fondamentales et Mathématiques Appliquées, ainsi que des masters en mathématiques.
    \item \textbf{Département d'Informatique :} responsable de la filière Licence MIP (Mathématiques, Informatique et Physique) et de plusieurs masters en informatique.
    \item \textbf{Département de Physique :} offre des licences en Physique Fondamentale et Physique des Matériaux.
    \item \textbf{Département de Chimie :} gère les licences en Chimie Fondamentale et Chimie Analytique.
    \item \textbf{Département de Biologie :} propose des licences en Biologie Moléculaire et Biotechnologie.
    \item \textbf{Département de Géologie :} offre des formations en Sciences de la Terre.
\end{itemize}

Chaque département est dirigé par un chef de département élu, assisté d'un conseil de département. Il coordonne les activités pédagogiques, les emplois du temps et la répartition des enseignements.

\subsection{Environnement du stage}
\label{ch1:environnement}

\subsubsection{Département d'accueil}
\label{ch1:equipe}

Notre projet de fin d'études s'est déroulé au sein de la filière \textbf{Licence Fondamentale MIP (Mathématiques, Informatique et Physique)} du Département d'Informatique. Cette filière vise à former des étudiants polyvalents possédant une solide base en mathématiques et une maîtrise des technologies de l'information. La formation couvre des domaines variés tels que le développement logiciel, les bases de données, les réseaux, l'algorithmique avancée et l'intelligence artificielle.

Le corps enseignant de la filière est composé de professeurs de l'enseignement supérieur, de professeurs habilités et de professeurs assistants, couvrant l'ensemble des spécialités du programme.

\subsubsection{Contexte du projet au sein de l'organisme}
\label{ch1:contexte_projet}

Le projet de système de gestion des soutenances et des jurys s'inscrit dans une volonté de modernisation des processus administratifs et pédagogiques au sein de la faculté. Actuellement, l'organisation des soutenances -- qu'il s'agisse des soutenances de PFE, de mémoires de master ou de projets de licence -- repose sur des méthodes manuelles et des outils bureautiques génériques.

Le département MIP, confronté à un volume croissant d'étudiants et à une complexité accrue dans la coordination des jurys, a identifié le besoin d'un outil informatique dédié. Ce projet a donc été proposé par \textbf{Pr. Mohammed AIT DAOUD} dans le cadre des projets de fin d'études de la licence MIP, avec pour objectif de fournir une solution concrète et opérationnelle répondant aux besoins réels du département.

\subsection{Conclusion}
\label{ch1:conclusion}

La Faculté des Sciences Ben M'Sick constitue un cadre académique riche et diversifié, offrant des formations scientifiques de qualité dans un large éventail de disciplines. Le département MIP, par sa proximité avec les technologies de l'information, est naturellement le lieu idéal pour le développement d'une solution numérique de gestion des soutenances. Ce contexte a directement inspiré les choix technologiques et architecturaux qui seront présentés dans les chapitres suivants.

% ═══════════════════════════════════════════════════════════════════════════════
%  CHAPITRE 2 : ÉTUDE DE L'EXISTANT ET TECHNOLOGIES UTILISÉES
% ═══════════════════════════════════════════════════════════════════════════════
\newpage
\section{Étude de l'existant et technologies utilisées}
\label{ch2}

\subsection{Introduction}
\label{ch2:intro}

Ce chapitre présente une analyse approfondie du système actuel d'organisation des soutenances au sein de la faculté, identifie ses lacunes et compare les solutions existantes sur le marché. Il détaille ensuite les choix technologiques retenus pour le développement de notre application en justifiant chaque décision. Enfin, l'architecture générale du système est présentée.

\subsection{Analyse de l'existant}
\label{ch2:existant}

\subsubsection{Description du système actuel}
\label{ch2:systeme_actuel}

L'organisation actuelle des soutenances repose sur une combinaison d'outils bureautiques génériques et de processus manuels :

\begin{itemize}[label=\textcolor{primary}{\textbullet}]
    \item \textbf{Outils :} utilisation intensive de tableurs \textbf{Excel} pour le suivi des étudiants et des plannings, et de \textbf{documents papier} pour les évaluations.
    \item \textbf{Communication :} échanges fragmentés par \textbf{courrier électronique} pour la collecte des disponibilités et l'envoi des convocations.
    \item \textbf{Logistique :} gestion indépendante des salles via des services tiers, souvent sans visibilité directe sur les occupations.
    \item \textbf{Évaluation :} saisie manuelle des notes sur fiches imprimées, suivie d'une re-saisie numérique pour le calcul des moyennes.
\end{itemize}

\subsubsection{Processus actuel détaillé}
\label{ch2:processus_actuel}

L'organisation d'une session de soutenances suit un cheminement complexe divisé en deux grandes phases opérationnelles.

\begin{figure}[H]
    \centering
    \pumlsvg[0.6\textwidth]{process_planification}
    \caption{Processus actuel -- Phase~1 : Planification}
    \label{fig:processus_actuel_phase1}
\end{figure}

La Figure \ref{fig:processus_actuel_phase1} illustre la nature intrinsèquement itérative de la phase de planification. Elle repose sur une communication asynchrone par e-mail qui génère fréquemment des goulots d'étranglement : la centralisation manuelle des informations (étudiants, sujets, encadrants) et la négociation des créneaux imposent de multiples allers-retours avant d'aboutir à un planning acceptable.

\begin{figure}[H]
    \centering
    \pumlsvg[0.6\textwidth]{process_execution}
    \caption{Processus actuel -- Phase~2 : Exécution}
    \label{fig:processus_actuel_phase2}
\end{figure}

La Figure \ref{fig:processus_actuel_phase2} détaille la transition vers la phase d'exécution. Cette étape est critique car elle comporte une vérification manuelle des conflits (double affectation d'un enseignant ou d'une salle) avant la diffusion officielle. Le cycle se termine par une collecte physique des notes et un archivage souffrant d'un manque de centralisation numérique.

\subsubsection{Analyse critique}
\label{ch2:analyse_critique}

\textbf{Avantages.} Le système actuel présente certains avantages dans des contextes de très petite taille : flexibilité d'adaptation à des cas particuliers, absence de dépendance technologique et charge de travail acceptable pour un volume très limité de soutenances.

\textbf{Inconvénients.} Les inconvénients deviennent prépondérants dès que le volume augmente. Les erreurs de planification sont fréquentes et souvent détectées tardivement, provoquant des annulations de dernière minute. La traçabilité des décisions est faible, la charge du coordinateur est considérable en fin de semestre, et la ressaisie répétée des données multiplie les risques d'erreurs de transcription.

\subsection{Étude des solutions similaires}
\label{ch2:solutions_similaires}

Plusieurs solutions logicielles existent sur le marché ou sous forme de projets open source. Les Tableaux \ref{tab:erp}, \ref{tab:planification_generique} et \ref{tab:opensource} présentent une analyse comparative qui a guidé le positionnement de notre projet.

\subsubsection{Systèmes universitaires généraux (ERP)}
\label{ch2:erp}

\begin{table}[H]
    \centering
    \small
    \begin{tabular}{lp{5.5cm}p{5.5cm}}
        \toprule
        \rowcolor{primary!10}
        \textbf{Solution} & \textbf{Description} & \textbf{Adéquation} \\
        \midrule
        \textbf{Apogée (Amue)} &
        Système de gestion des étudiants utilisé dans les universités françaises &
        Trop généraliste ; gestion des soutenances non native \\
        \addlinespace
        \textbf{Aurion} &
        ERP académique couvrant scolarité, stages et PFE &
        Module soutenance limité, coût élevé, complexité d'intégration \\
        \addlinespace
        \textbf{Koha / Aleph} &
        Systèmes bibliothèques universitaires &
        Hors sujet \\
        \bottomrule
    \end{tabular}
    \caption{Analyse des systèmes ERP académiques}
    \label{tab:erp}
\end{table}

\subsubsection{Outils de planification générique}
\label{ch2:planification_generique}

\begin{table}[H]
    \centering
    \small
    \begin{tabular}{lp{5.5cm}p{5.5cm}}
        \toprule
        \rowcolor{primary!10}
        \textbf{Solution} & \textbf{Description} & \textbf{Limites pour ce besoin} \\
        \midrule
        \textbf{Doodle}          & Sondages de disponibilité & Pas de gestion des jurys ni de génération PDF \\
        \addlinespace
        \textbf{Google Calendar} & Calendrier partagé        & Pas de logique métier académique, pas de convocations \\
        \addlinespace
        \textbf{Microsoft Teams} & Gestion de réunions       & Non adapté à la logique jury/soutenance \\
        \bottomrule
    \end{tabular}
    \caption{Analyse des outils de planification génériques}
    \label{tab:planification_generique}
\end{table}

\subsubsection{Solutions open source (GitHub)}
\label{ch2:opensource}

\begin{table}[H]
    \centering
    \small
    \begin{tabular}{lp{5cm}p{6cm}}
        \toprule
        \rowcolor{primary!10}
        \textbf{Projet}          & \textbf{Lien GitHub} & \textbf{Limites identifiées} \\
        \midrule
        \textbf{appliSoutenance} & \href{https://github.com/cdegrel/appliSoutenance}{cdegrel/appliSoutenance} & Focalisé sur la notation en séance ; ne gère pas la planification amont ni les salles \\
        \addlinespace
        \textbf{PFE-Management}  & \href{https://github.com/khaireddine24/PFE-Management-System}{khaireddine24/PFE-MS} & Gestion administrative de base ; absence de détection de conflits et de génération PDF \\
        \addlinespace
        \textbf{Jury (HackUTD)}  & \href{https://github.com/hackutd/jury}{hackutd/jury} & Adapté aux hackathons ; incompatible avec les rôles académiques (Président, Rapporteur) \\
        \bottomrule
    \end{tabular}
    \caption{Analyse de projets open source similaires}
    \label{tab:opensource}
\end{table}

\begin{projectbox}{Conclusion de l'analyse}
    Comme le montrent les Tableaux \ref{tab:erp}, \ref{tab:planification_generique} et \ref{tab:opensource}, aucune solution existante ne réunit simultanément les quatre points suivants :
    \begin{itemize}[label=\textcolor{primary}{\textbullet}, leftmargin=1cm]
        \item Gestion complète du cycle de soutenance (planification \(\rightarrow\) évaluation).
        \item Composition et affectation de jurys avec détection de conflits.
        \item Génération automatique de convocations et fiches d'évaluation.
        \item Adaptation au contexte marocain (filières, CNE, structure universitaire locale).
    \end{itemize}
\end{projectbox}

\subsection{Choix technologiques}
\label{ch2:choix_technologiques}

Le choix des technologies est une étape cruciale qui conditionne la réussite du projet. Chaque décision a été prise en fonction des besoins spécifiques du système, des compétences de l'équipe et des contraintes du projet.

\subsubsection{Frontend}
\label{ch2:frontend}

\textbf{React.js} a été retenu pour le développement de l'interface utilisateur. Ce choix se justifie par :
\begin{itemize}[label=\textcolor{primary}{\textbullet}]
    \item \textbf{Réactivité :} son modèle de composants et le DOM virtuel offrent une expérience utilisateur fluide et sans rechargement de page.
    \item \textbf{Écosystème riche :} une vaste communauté, des bibliothèques matures et des outils de développement performants.
    \item \textbf{Compatibilité REST :} s'intègre parfaitement avec une API REST Spring Boot.
\end{itemize}
L'interface est conçue avec \textbf{Tailwind CSS} et \textbf{Shadcn UI} pour un design responsive, moderne et cohérent.

\subsubsection{Backend}
\label{ch2:backend}

\textbf{Spring Boot} (Java) a été choisi pour la couche serveur pour les raisons suivantes :
\begin{itemize}[label=\textcolor{primary}{\textbullet}]
    \item \textbf{Robustesse :} framework éprouvé pour les applications d'entreprise, offrant une sécurité et une stabilité élevées.
    \item \textbf{Spring Security :} intégration native pour la gestion de l'authentification JWT et le contrôle d'accès par rôles.
    \item \textbf{Spring Data JPA :} facilite l'accès aux données et la gestion des entités relationnelles.
    \item \textbf{Architecture MVC :} séparation claire des préoccupations (Contrôleurs, Services, Repositories).
    \item \textbf{Documentation et communauté :} large base de connaissances et support actif.
\end{itemize}

\subsubsection{Base de données}
\label{ch2:base_donnees}

\textbf{MySQL} a été sélectionné comme système de gestion de base de données :
\begin{itemize}[label=\textcolor{primary}{\textbullet}]
    \item \textbf{Relationnel :} parfaitement adapté à la structure de données fortement interconnectée du système (étudiants, projets, jurys, soutenances, évaluations).
    \item \textbf{Performance :} fiable et performant pour les charges de travail académiques.
    \item \textbf{Intégration :} compatibilité native avec Spring Data JPA et Hibernate.
    \item \textbf{Coût :} solution open source mature sans frais de licence.
\end{itemize}

\subsubsection{Outils de développement}
\label{ch2:outils_dev}

\begin{itemize}[label=\textcolor{primary}{\textbullet}]
    \item \textbf{Git :} versionnement du code source et collaboration via GitHub.
    \item \textbf{VS Code :} environnement de développement intégré pour le frontend.
    \item \textbf{IntelliJ IDEA :} IDE dédié au développement Spring Boot.
    \item \textbf{Postman :} test et documentation des API REST.
    \item \textbf{Maven :} gestion des dépendances et construction du projet backend.
\end{itemize}

\subsection{Architecture générale retenue}
\label{ch2:architecture}

L'application adopte une architecture web en trois couches classique (Figure \ref{fig:architecture_generale}), favorisant la modularité et la séparation des préoccupations :

\begin{itemize}[label=\textcolor{primary}{\textbullet}]
    \item \textbf{Couche Frontend :} interface utilisateur React, affichage des données et capture des interactions.
    \item \textbf{Couche Backend :} logique métier Spring Boot, règles de gestion, détection des conflits et génération des documents.
    \item \textbf{Couche Base de données :} persistance MySQL, intégrité et sécurité des informations du système.
\end{itemize}

\begin{figure}[H]
    \centering
    \pumlsvg[0.9\textwidth]{arch_layers}
    \caption{Architecture trois couches du système}
    \label{fig:architecture_generale}
\end{figure}

\subsection{Conclusion}
\label{ch2:conclusion}

Ce chapitre a dressé un état des lieux complet du système actuel d'organisation des soutenances, mettant en évidence ses limites opérationnelles. L'analyse des solutions existantes a confirmé l'absence d'un outil répondant spécifiquement aux besoins de notre contexte. Les choix technologiques -- React, Spring Boot, MySQL -- ont été justifiés et l'architecture à trois couches a été présentée comme fondation du système. Le chapitre suivant détaillera l'analyse des besoins fonctionnels et la conception complète du système.

% ═══════════════════════════════════════════════════════════════════════════════
%  CHAPITRE 3 : ANALYSE DES BESOINS ET CONCEPTION DU SYSTÈME
% ═══════════════════════════════════════════════════════════════════════════════
\newpage
\section{Analyse des besoins et conception du système}
\label{ch3}

\subsection{Introduction}
\label{ch3:intro}

Ce chapitre constitue le cœur de la phase d'analyse et de conception du système. Il définit le cahier des charges fonctionnel, présente les règles de gestion qui encadrent le système, détaille la modélisation UML complète -- des cas d'utilisation aux diagrammes de séquence --, décrit l'architecture technique et le modèle de données relationnel, et enfin présente la conception de l'interface utilisateur.

\subsection{Présentation générale du projet}
\label{ch3:presentation_generale}

Le projet consiste à développer une application web pour la gestion des soutenances et des jurys au sein de la Faculté des Sciences Ben M'Sick. Le système permet de :
\begin{itemize}[label=\textcolor{primary}{\textbullet}]
    \item Planifier les soutenances individuelles au sein de sessions globales définies par l'administration.
    \item Affecter automatiquement les jurys en respectant les contraintes de charge et de spécialité.
    \item Détecter les conflits horaires, de salle et de double affectation.
    \item Générer les convocations et documents administratifs.
    \item Gérer les évaluations et produire les résultats consolidés.
\end{itemize}

\subsection{Cahier des charges}
\label{ch3:cahier_charges}

\subsubsection{Acteurs du système}
\label{ch3:acteurs}

Le système doit servir quatre catégories d'acteurs distinctes, chacune disposant de droits d'accès et de fonctionnalités spécifiques :

\begin{itemize}[label=\textcolor{primary}{\textbullet}]
    \item \textbf{Administrateur :} gère les comptes utilisateurs, configure les paramètres généraux et supervise le bon fonctionnement du système.
    \item \textbf{Coordinateur :} acteur central de l'organisation ; planifie les soutenances dans une session globale, affecte les jurys, valide les plannings et consulte les rapports de synthèse.
    \item \textbf{Enseignant / Membre de jury :} consulte ses convocations, prend connaissance des informations relatives aux étudiants à évaluer et saisit les notes et observations après chaque soutenance.
    \item \textbf{Étudiant :} consulte la date, l'heure, la salle et la composition du jury pour sa soutenance, et peut télécharger sa convocation.
\end{itemize}

\subsubsection{Besoins fonctionnels}
\label{ch3:besoins_fonctionnels}

\textbf{Gestion des utilisateurs.} Le système doit permettre la création et la gestion des comptes utilisateurs avec une authentification sécurisée. Les quatre rôles définis -- administrateur, coordinateur, enseignant et étudiant -- doivent chacun disposer d'un ensemble de droits spécifiques.

\textbf{Gestion des étudiants.} L'application doit permettre l'enregistrement des étudiants avec leurs informations complètes : nom, prénom, code national d'étudiant (CNE), filière, niveau, titre du projet et encadrant. Une fonction d'import depuis un fichier CSV ou Excel est prévue. Chaque étudiant doit être associé à un projet spécifique.

\textbf{Gestion des enseignants et jurys.} Le système doit permettre l'ajout et la gestion des enseignants. Par défaut, tous les enseignants sont considérés comme disponibles pour les jurys, mais ils peuvent déclarer leurs plages d'indisponibilité. La composition d'un jury doit inclure la désignation d'un président, d'un rapporteur, d'un examinateur et, le cas échéant, de l'encadrant de l'étudiant.

\textbf{Gestion des salles.} La gestion des salles comprend l'ajout de nouvelles salles, le renseignement de leur capacité, la consultation de leur disponibilité et l'affectation à une soutenance spécifique.

\textbf{Planification des soutenances.} La planification permet la création des soutenances individuelles au sein d'une session globale définie par l'administration, avec la date, l'heure de début, la durée et la salle. Chaque soutenance doit être associée à un étudiant ou à un groupe, et un planning global doit être consultable selon différents critères.

\textbf{Affectation des jurys.} L'affectation des jurys consiste à associer un ensemble d'enseignants à chaque soutenance. Le système doit vérifier qu'un enseignant n'est pas affecté à deux soutenances simultanées et doit rechercher une répartition équilibrée de la charge.

\textbf{Gestion des conflits.} La détection des conflits est une fonctionnalité critique. Le système doit détecter et signaler -- avant validation -- les conflits de salle, les conflits d'horaire pour un enseignant, la double planification d'un même étudiant et les doublons de planification.

\textbf{Génération de documents.} Le système doit générer : convocations des étudiants, convocations des jurys, planning général, fiches d'évaluation, listes de présence et procès-verbal simplifié. Les formats supportés sont le PDF, l'impression directe et l'export Excel/CSV.

\textbf{Saisie des évaluations.} Les membres du jury ou le coordinateur doivent pouvoir saisir les notes individuelles, ajouter des remarques, puis consolider le résultat final. La décision finale (admis, admis avec corrections, ajourné) est enregistrée après calcul de la moyenne.

\textbf{Tableau de bord.} Un tableau de bord synthétique doit afficher le nombre total de soutenances, le nombre de jurys mobilisés, les salles utilisées, les soutenances planifiées par jour ainsi que le suivi des soutenances évaluées et restant à évaluer.

\subsubsection{Besoins non fonctionnels}
\label{ch3:besoins_non_fonctionnels}

\textbf{Performance.} L'affichage d'un planning ne doit pas excéder deux secondes. La recherche d'un utilisateur doit être quasi instantanée. La génération d'un document PDF ne doit pas dépasser cinq secondes dans des conditions normales d'utilisation.

\textbf{Sécurité.} L'accès est protégé par une authentification par identifiant et mot de passe. Les mots de passe sont stockés sous forme de hachage cryptographique (BCrypt). Les fonctionnalités et les données sont strictement restreintes selon le rôle de l'utilisateur connecté.

\textbf{Fiabilité.} Le système doit éviter toute perte de données lors des opérations de création ou de modification. Les erreurs de saisie doivent être signalées par des messages compréhensibles. En cas de défaillance, le système doit pouvoir restaurer son état antérieur sans corruption des données.

\textbf{Maintenabilité.} Le code source doit être structuré, documenté et versionné avec Git. Une architecture modulaire facilitera les évolutions ultérieures.

\subsection{Règles de gestion}
\label{ch3:regles_gestion}

\begin{projectbox}{Règles de gestion (RG)}
    \begin{itemize}[label=\textcolor{primary}{\textbullet}, leftmargin=1.5cm]
        \item \textbf{RG1 :} Un enseignant ne peut être affecté à deux soutenances ayant des horaires qui se chevauchent.
        \item \textbf{RG2 :} Une salle ne peut être réservée pour deux soutenances simultanées.
        \item \textbf{RG3 :} Un étudiant ne peut avoir qu'une seule soutenance programmée.
        \item \textbf{RG4 :} Un jury doit comporter au minimum trois membres (Président, Rapporteur, Examinateur). L'encadrant de l'étudiant peut également en faire partie.
        \item \textbf{RG5 :} L'encadrant de l'étudiant ne peut en aucun cas occuper le rôle de président du jury.
        \item \textbf{RG6 :} La note finale est calculée comme la moyenne des notes individuelles validées par les membres du jury.
    \end{itemize}
\end{projectbox}

Ces règles sont appliquées automatiquement lors de chaque opération de planification ou d'affectation.

\subsection{Modules du système}
\label{ch3:modules}

Dans ce document, le terme \textit{module} désigne une unité fonctionnelle du système, regroupant un ensemble cohérent de fonctionnalités liées à un domaine métier spécifique. L'application est structurée en sept modules fonctionnels couvrant l'ensemble du cycle de gestion des soutenances.

\subsubsection{Module 1 : Authentification et gestion des rôles}
\label{ch3:module_auth}

Ce module gère l'accès sécurisé à la plateforme. Il assure la connexion des utilisateurs, la gestion de leurs profils et l'attribution des droits selon les quatre rôles définis.

\subsubsection{Module 2 : Gestion académique}
\label{ch3:module_academique}

Ce module centralise les informations relatives aux étudiants, aux projets de fin d'études, aux encadrants et aux enseignants.

\subsubsection{Module 3 : Planification}
\label{ch3:module_planification}

Ce module permet de créer et d'organiser les soutenances individuelles au sein d'une session globale : définition des horaires, réservation des salles et visualisation du planning global.

\subsubsection{Module 4 : Gestion des jurys}
\label{ch3:module_jurys}

Ce module facilite l'affectation des membres de jury à chaque soutenance, la composition équilibrée des commissions, ainsi que la gestion et la consultation des indisponibilités déclarées par les enseignants.

\subsubsection{Module 5 : Évaluation}
\label{ch3:module_evaluation}

Dédié au processus d'évaluation, ce module permet la saisie des notes, l'ajout d'observations, la génération des fiches d'évaluation et l'enregistrement des décisions de délibération.

\subsubsection{Module 6 : Génération documentaire}
\label{ch3:module_documents}

Ce module automatise la production des documents administratifs essentiels : convocations, plannings, procès-verbaux et exports de données.

\subsubsection{Module 7 : Tableau de bord}
\label{ch3:module_dashboard}

Le tableau de bord fournit une vue synthétique de l'activité : statistiques sur les soutenances, synthèse des évaluations et suivi en temps réel de l'état d'avancement des sessions.

\subsection{Conception UML}
\label{ch3:conception_uml}

Cette section détaille la modélisation UML complète du système. Les diagrammes présentés couvrent l'ensemble des vues nécessaires à la compréhension du système : vue fonctionnelle (cas d'utilisation), vue statique (classes), vue dynamique (séquences, activités, états).

\subsubsection{Diagrammes de cas d'utilisation}
\label{ch3:uc}

Les diagrammes de cas d'utilisation permettent de délimiter le périmètre du projet en identifiant les interactions entre les différents acteurs et les fonctionnalités offertes par l'application. Chaque module répond à des objectifs spécifiques, garantissant une séparation claire des responsabilités.

\begin{projectbox}{Remarque}
    Par souci de clarté visuelle, le cas d'utilisation « S'authentifier » est omis des diagrammes détaillés mais reste une pré-condition transversale à toutes les actions du système.
\end{projectbox}

\paragraph{Hiérarchie des acteurs.} Le système distingue plusieurs profils d'utilisateurs, chacun possédant des droits et des responsabilités spécifiques. La hiérarchie permet d'hériter des fonctionnalités communes, notamment la consultation du planning et la gestion du profil personnel.

\begin{figure}[H]
    \centering
    \pumlsvg[0.6\textwidth]{uc_hierarchy}
    \caption{Hiérarchie générale des acteurs}
    \label{fig:uc_hierarchy}
\end{figure}

\paragraph{Module Administration.} Ce module est réservé à la gestion technique et structurelle de la plateforme. L'administrateur veille au bon fonctionnement du système en gérant les ressources de base.

\begin{figure}[H]
    \centering
    \pumlsvg[0.7\textwidth]{uc_admin}
    \caption{Cas d'utilisation : Administration du système}
    \label{fig:uc_admin}
\end{figure}

\begin{table}[H]
    \centering
    \small
    \begin{tabular}{lp{10cm}}
        \toprule
        \rowcolor{primary!10}
        \textbf{Cas d'utilisation} & \textbf{Description} \\
        \midrule
        Gérer comptes utilisateurs & Création, modification (incluant rôles et mots de passe) et suppression des comptes. \\
        \addlinespace
        Gérer salles & Administration du catalogue des salles, capacité et disponibilités. \\
        \addlinespace
        Gérer sessions globales & Création, modification et clôture des sessions globales (normale, rattrapage). \\
        \addlinespace
        Gérer paramètres système & Configuration des paramètres transversaux : durées par défaut, périodes, règles générales. \\
        \bottomrule
    \end{tabular}
    \caption{Description des cas d'utilisation -- Administrateur}
    \label{tab:uc_admin}
\end{table}

\paragraph{Module Coordination.} Le coordinateur est l'acteur central du processus de planification. Dans ce module, le terme \textit{soutenance} désigne une séance individuelle planifiée pour un étudiant et son jury, tandis que la \textit{session globale} désigne le cadre administratif créé par l'administrateur. Le coordinateur travaille sur les soutenances contenues dans une session globale, de l'importation massive des données à la validation finale du planning.

\begin{figure}[H]
    \centering
    \pumlsvg[\textwidth]{uc_coordination}
    \caption{Cas d'utilisation : Coordination des soutenances}
    \label{fig:uc_coordination}
\end{figure}

\begin{table}[H]
    \centering
    \small
    \begin{tabular}{lp{10cm}}
        \toprule
        \rowcolor{primary!10}
        \textbf{Cas d'utilisation} & \textbf{Description} \\
        \midrule
        Assigner jurys & Constitution des commissions (président, rapporteur, examinateur). \\
        \addlinespace
        Planifier soutenances & Planification des créneaux, définition des horaires et affectation des salles. \\
        \addlinespace
        Valider planning & Confirmation finale du planning après traitement des alertes et conflits. \\
        \addlinespace
        Gérer projets & Gestion des thèmes, titres de PFE et association étudiants-projets-encadrants. \\
        \addlinespace
        Consulter planning & Visualisation globale filtrable (date, salle, enseignant, filière, étudiant). \\
        \addlinespace
        Rechercher / Filtrer & Recherche multicritère avancée pour la supervision. \\
        \addlinespace
        Générer convocations & Production automatique des documents PDF pour étudiants et jurys. \\
        \addlinespace
        Générer planning général & Production du planning complet et liste de présence pour la session. \\
        \addlinespace
        Consulter rapports & Accès aux synthèses, statistiques et procès-verbaux simplifiés. \\
        \addlinespace
        Vérifier conflits & Détection des chevauchements, doubles planifications et doublons. \\
        \bottomrule
    \end{tabular}
    \caption{Description des cas d'utilisation -- Coordinateur}
    \label{tab:uc_coordination}
\end{table}

\paragraph{Module Jury.} Les enseignants interagissent avec le système principalement pour consulter leurs obligations de jury et évaluer les travaux des étudiants.

\begin{figure}[H]
    \centering
    \pumlsvg[\textwidth]{uc_jury}
    \caption{Cas d'utilisation : Enseignant et Jury}
    \label{fig:uc_jury}
\end{figure}

\begin{table}[H]
    \centering
    \small
    \begin{tabular}{lp{10cm}}
        \toprule
        \rowcolor{primary!10}
        \textbf{Cas d'utilisation} & \textbf{Description} \\
        \midrule
        Consulter planning & Accès au planning personnel des soutenances. \\
        \addlinespace
        Consulter convocation & Visualisation et téléchargement des convocations. \\
        \addlinespace
        Saisir évaluation & Enregistrement des notes et observations pendant ou après la soutenance. \\
        \addlinespace
        Valider / modifier évaluation & Finalisation de la notation et modification des remarques. \\
        \addlinespace
        Consulter informations étudiant & Consultation des détails du projet et du dossier. \\
        \addlinespace
        Générer fiches d'évaluation & Création des supports de notation pour la session. \\
        \addlinespace
        Déclarer indisponibilité & Signalement des créneaux horaires non disponibles. \\
        \bottomrule
    \end{tabular}
    \caption{Description des cas d'utilisation -- Enseignant}
    \label{tab:uc_jury}
\end{table}

\paragraph{Module Étudiant.} L'interface dédiée à l'étudiant est conçue pour être simple et directe. Elle lui permet de se concentrer sur les informations essentielles liées à sa soutenance.

\begin{figure}[H]
    \centering
    \pumlsvg[\textwidth]{uc_student}
    \caption{Cas d'utilisation : Espace Étudiant}
    \label{fig:uc_student}
\end{figure}

\begin{table}[H]
    \centering
    \small
    \begin{tabular}{lp{10cm}}
        \toprule
        \rowcolor{primary!10}
        \textbf{Cas d'utilisation} & \textbf{Description} \\
        \midrule
        Télécharger convocation & Récupération du document PDF officiel. \\
        \addlinespace
        Consulter infos de soutenance & Visualisation de la date, l'heure, la salle et le jury. \\
        \addlinespace
        Consulter résultats & Consultation des notes et de la décision finale. \\
        \bottomrule
    \end{tabular}
    \caption{Description des cas d'utilisation -- Étudiant}
    \label{tab:uc_student}
\end{table}

\subsubsection{Diagramme de classes}
\label{ch3:classes}

Le diagramme de classes constitue le pivot de la conception statique. Il définit la structure de données interne et les règles métier qui régissent les interactions entre les objets du système.

\begin{figure}[H]
    \centering
    \pumlsvg[\textwidth]{class_global}
    \caption{Diagramme de classes global du système}
    \label{fig:class_global}
\end{figure}

Les entités clés identifiées sont les suivantes :
\begin{itemize}[label=\textcolor{primary}{\textbullet}]
    \item \textbf{Utilisateur :} Classe de base gérant l'identité, les identifiants de connexion et les droits d'accès (rôles).
    \item \textbf{Soutenance :} Entité centrale faisant le lien entre un projet, un jury, une salle et un créneau temporel au sein d'une session globale.
    \item \textbf{Projet :} Représente le travail académique (PFE, Mémoire, Stage), supervisé par un encadrant.
    \item \textbf{Jury :} Structure d'association définissant les rôles (Président, Rapporteur, Examinateur).
    \item \textbf{Évaluation et résultat :} Chaque membre du jury saisit une note individuelle ; le résultat consolide ces notes en moyenne finale.
\end{itemize}

\subsubsection{Cycle de vie d'une soutenance}
\label{ch3:etat_soutenance}

Le diagramme d'états ci-dessous illustre les transitions de l'entité Soutenance. Ce cycle garantit que les actions (saisie des notes, génération des convocations) ne sont autorisées que lorsque la soutenance est dans l'état approprié.

\begin{figure}[H]
    \centering
    \pumlsvg[0.6\textwidth]{state_soutenance}
    \caption{Diagramme d'états : Cycle de vie d'une soutenance}
    \label{fig:state_soutenance}
\end{figure}

\subsubsection{Diagrammes d'activité}
\label{ch3:activite}

Les diagrammes d'activité détaillent la logique algorithmique et le flux de contrôle des processus métier critiques.

\textbf{Workflow de planification.} Ce processus décrit les étapes de planification des soutenances à l'intérieur d'une session globale. Le workflow est itératif, permettant au coordinateur de corriger les anomalies détectées (conflits de salle, indisponibilité) avant la validation finale.

\begin{figure}[H]
    \centering
    \pumlsvg[0.6\textwidth]{act_planning}
    \caption{Diagramme d'activité : Planification de session}
    \label{fig:act_planning}
\end{figure}

\textbf{Déroulement des évaluations.} Ce diagramme illustre le processus d'évaluation, du contrôle de présence jusqu'au calcul automatique de la moyenne finale.

\begin{figure}[H]
    \centering
    \pumlsvg[0.58\textwidth]{act_evaluation}
    \caption{Diagramme d'activité : Processus d'évaluation}
    \label{fig:act_evaluation}
\end{figure}

\subsubsection{Diagrammes de séquence}
\label{ch3:sequence}

Les diagrammes de séquence détaillent les interactions chronologiques entre les composants (Client, Contrôleur, Services, Modèle) lors de l'exécution des processus métier.

\textbf{Processus d'authentification.} L'authentification repose sur un échange sécurisé où le contrôleur valide les identifiants via le service utilisateur avant de générer un token JWT.

\begin{figure}[H]
    \centering
    \pumlsvg[\textwidth]{seq_auth}
    \caption{Diagramme de séquence : Authentification}
    \label{fig:seq_auth}
\end{figure}

\textbf{Planification et détection de conflits.} Lors de la planification, le PlanningService sollicite le ConflictService pour vérifier l'absence de chevauchements.

\begin{figure}[H]
    \centering
    \pumlsvg[\textwidth]{seq_planning}
    \caption{Diagramme de séquence : Planification de soutenance}
    \label{fig:seq_planning}
\end{figure}

\textbf{Affectation et validation du jury.} L'affectation des jurys illustre les règles métier RG4 et RG5. Le service métier coordonne la sélection des enseignants selon la charge de travail et la spécialité.

\begin{figure}[H]
    \centering
    \pumlsvg[\textwidth]{seq_jury_assignment}
    \caption{Diagramme de séquence : Affectation du jury}
    \label{fig:seq_jury}
\end{figure}

\textbf{Saisie et validation d'une évaluation.} Ce flux montre comment les membres du jury enregistrent leurs notes. Le système recalcule la moyenne provisoire à chaque saisie et ne finalise le résultat que lorsque toutes les notes sont validées.

\begin{figure}[H]
    \centering
    \pumlsvg[\textwidth]{seq_evaluation}
    \caption{Diagramme de séquence : Saisie d'évaluation}
    \label{fig:seq_evaluation}
\end{figure}

\textbf{Génération PDF.} Le service de documents extrait les données consolidées du modèle pour les injecter dans un template normalisé.

\begin{figure}[H]
    \centering
    \pumlsvg[\textwidth]{seq_pdf_generation}
    \caption{Diagramme de séquence : Génération PDF}
    \label{fig:seq_pdf}
\end{figure}

\textbf{Import massif de données.} Le processus d'importation gère la lecture de fichiers externes, effectue une validation ligne par ligne et assure la persistance transactionnelle.

\begin{figure}[H]
    \centering
    \pumlsvg[\textwidth]{seq_import_data}
    \caption{Diagramme de séquence : Import massif de données}
    \label{fig:seq_import}
\end{figure}

\textbf{Validation finale et notifications.} La validation par le coordinateur verrouille le planning et déclenche des événements asynchrones (génération et envoi des notifications).

\begin{figure}[H]
    \centering
    \pumlsvg[\textwidth]{seq_validate_planning}
    \caption{Diagramme de séquence : Validation finale et notifications}
    \label{fig:seq_validation}
\end{figure}

\textbf{Récupération de compte.} En cas d'oubli du mot de passe, un lien temporaire contenant un jeton unique est envoyé pour une réinitialisation sécurisée.

\begin{figure}[H]
    \centering
    \pumlsvg[\textwidth]{seq_reset_password}
    \caption{Diagramme de séquence : Réinitialisation de mot de passe}
    \label{fig:seq_reset_password}
\end{figure}

\subsection{Architecture technique}
\label{ch3:architecture_technique}

Le système repose sur une architecture web moderne organisée en couches, garantissant une séparation claire des responsabilités, une sécurité renforcée et une maintenance facilitée.

\subsubsection{Architecture en couches}
\label{ch3:couches}

L'application suit le pattern \textbf{MVC (Modèle-Vue-Contrôleur)} adapté au Web :
\begin{itemize}[label=\textcolor{primary}{\textbullet}]
    \item \textbf{Vue (View) :} Assurée par le Frontend React, elle gère l'interface utilisateur et l'expérience interactive de manière asynchrone.
    \item \textbf{Contrôleur (Controller) :} Les contrôleurs Spring Boot exposent des endpoints REST, reçoivent les requêtes du client et orchestrent les réponses.
    \item \textbf{Modèle (Model) :} Comprend la logique métier (Services), l'accès aux données (Repositories) et la structure de persistance MySQL.
\end{itemize}

\begin{figure}[H]
    \centering
    \pumlsvg[0.8\textwidth]{arch_layers}
    \caption{Architecture générale en trois couches}
    \label{fig:arch_layers}
\end{figure}

\subsubsection{Architecture technique des composants}
\label{ch3:composants}

Le diagramme suivant détaille l'organisation interne des modules et leur mode de communication. On y observe la séparation nette entre le client web et le serveur d'application qui encapsule la logique métier complexe.

\begin{figure}[H]
    \centering
    \pumlsvg[\textwidth]{component_architecture}
    \caption{Architecture technique des composants}
    \label{fig:component_architecture}
\end{figure}

\subsubsection{Sécurité et Authentification JWT}
\label{ch3:securite_jwt}

Pour répondre aux exigences de sécurité et de performance, le système implémente une authentification \textbf{stateless} basée sur les \textbf{JSON Web Tokens (JWT)} :

\begin{itemize}[label=\textcolor{primary}{\textbullet}]
    \item \textbf{Génération :} Lors de la connexion, le serveur valide les identifiants et génère un token signé cryptographiquement contenant l'identité et le rôle de l'utilisateur.
    \item \textbf{Transmission :} Ce token est renvoyé au client (React) qui le stocke localement et l'inclut dans le header \texttt{Authorization} de chaque requête API ultérieure.
    \item \textbf{Validation :} Le backend (Spring Security) intercepte chaque requête, vérifie la signature du token et extrait les droits d'accès sans avoir à interroger la base de données pour la session.
\end{itemize}

\subsection{Modélisation de la base de données}
\label{ch3:mld}

Le Modèle Logique de Données (MLD) traduit les diagrammes de classes en structures tabulaires optimisées pour la persistance dans une base de données relationnelle MySQL.

Le modèle distingue explicitement les \textbf{sessions globales}, administrées au niveau institutionnel, et les \textbf{soutenances}, qui représentent les passages individuels planifiés par le coordinateur. Il ajoute également les tables de support nécessaires : indisponibilités des enseignants, convocations, notifications et jetons de réinitialisation.

\begin{figure}[H]
    \centering
    \pumlsvg[\textwidth]{mld_schema}
    \caption{Modèle Logique de Données (MLD)}
    \label{fig:mld_schema}
\end{figure}

\subsubsection{Dictionnaire des données}
\label{ch3:dictionnaire}

Ce dictionnaire détaille les attributs de l'ensemble des entités du MLD, organisés par domaine fonctionnel.

\paragraph{Gestion des utilisateurs et des rôles.}

\begin{table}[H]
    \centering
    \small
    \begin{tabular}{lp{2cm}lp{6cm}}
        \toprule
        \rowcolor{primary!10}
        \textbf{Champ} & \textbf{Type} & \textbf{Contrainte} & \textbf{Description} \\
        \midrule
        id             & INT           & PK, AI              & Identifiant unique de l'utilisateur. \\
        nom, prenom    & VARCHAR       & NOT NULL            & Informations civiles. \\
        email          & VARCHAR       & UNIQUE              & Adresse de contact et identifiant de connexion. \\
        password       & VARCHAR       & NOT NULL            & Hash du mot de passe (BCrypt). \\
        role           & ENUM          & NOT NULL            & Rôle (ADMIN, COORDINATEUR, ENSEIGNANT, ETUDIANT). \\
        actif          & BOOLEAN       & DEFAULT TRUE        & État d'activation du compte. \\
        \bottomrule
    \end{tabular}
    \caption{Dictionnaire -- Table utilisateurs}
    \label{tab:table_utilisateurs}
\end{table}

\begin{table}[H]
    \centering
    \small
    \begin{tabular}{lp{2cm}lp{6cm}}
        \toprule
        \rowcolor{primary!10}
        \textbf{Champ} & \textbf{Type} & \textbf{Contrainte} & \textbf{Description} \\
        \midrule
        cne            & VARCHAR       & PK                  & Code National de l'Étudiant. \\
        filiere        & VARCHAR       & NOT NULL            & Nom de la filière d'attachement. \\
        niveau         & VARCHAR       & NOT NULL            & Niveau académique (ex: Licence, Master). \\
        user\_id       & INT           & FK                  & Référence vers le compte utilisateur. \\
        projet\_id     & INT           & FK                  & Référence vers le projet porté par l'étudiant. \\
        \bottomrule
    \end{tabular}
    \caption{Dictionnaire -- Table etudiants}
    \label{tab:table_etudiants}
\end{table}

\begin{table}[H]
    \centering
    \small
    \begin{tabular}{lp{2cm}lp{6cm}}
        \toprule
        \rowcolor{primary!10}
        \textbf{Champ} & \textbf{Type} & \textbf{Contrainte} & \textbf{Description} \\
        \midrule
        id             & INT           & PK, AI              & Identifiant unique de l'enseignant. \\
        grade          & VARCHAR       & -                   & Grade académique (ex: PES, PH). \\
        specialite     & VARCHAR       & -                   & Domaine d'expertise. \\
        user\_id       & INT           & FK                  & Référence vers le compte utilisateur. \\
        \bottomrule
    \end{tabular}
    \caption{Dictionnaire -- Table enseignants}
    \label{tab:table_enseignants}
\end{table}

\begin{table}[H]
    \centering
    \small
    \begin{tabular}{lp{2cm}lp{6cm}}
        \toprule
        \rowcolor{primary!10}
        \textbf{Champ} & \textbf{Type} & \textbf{Contrainte} & \textbf{Description} \\
        \midrule
        id             & INT           & PK, AI              & Identifiant unique du profil spécifique. \\
        user\_id       & INT           & FK                  & Référence vers le compte utilisateur. \\
        \bottomrule
    \end{tabular}
    \caption{Dictionnaire -- Table administrateurs / coordinateurs}
    \label{tab:table_admin_coord}
\end{table}

\paragraph{Organisation académique et logistique.}

\begin{table}[H]
    \centering
    \small
    \begin{tabular}{lp{2cm}lp{6cm}}
        \toprule
        \rowcolor{primary!10}
        \textbf{Champ} & \textbf{Type} & \textbf{Contrainte} & \textbf{Description} \\
        \midrule
        id             & INT           & PK, AI              & Identifiant de la session globale. \\
        libelle        & VARCHAR       & NOT NULL            & Nom de la session (normale, rattrapage). \\
        date\_debut    & DATE          & NOT NULL            & Début de la période de soutenances. \\
        date\_fin      & DATE          & NOT NULL            & Fin de la période de soutenances. \\
        statut         & ENUM          & NOT NULL            & État (OUVERTE, VERROUILLEE, CLOTUREE). \\
        \bottomrule
    \end{tabular}
    \caption{Dictionnaire -- Table sessions\_globales}
    \label{tab:table_sessions}
\end{table}

\begin{table}[H]
    \centering
    \small
    \begin{tabular}{lp{2cm}lp{6cm}}
        \toprule
        \rowcolor{primary!10}
        \textbf{Champ} & \textbf{Type} & \textbf{Contrainte} & \textbf{Description} \\
        \midrule
        id             & INT           & PK, AI              & Identifiant unique du projet. \\
        titre          & VARCHAR       & NOT NULL            & Intitulé du sujet de PFE/Mémoire. \\
        description    & TEXT          & -                   & Résumé ou détails du projet. \\
        type           & ENUM          & NOT NULL            & Type (PFE, MEMOIRE, STAGE). \\
        encadrant\_id  & INT           & FK                  & Référence vers l'enseignant encadrant. \\
        \bottomrule
    \end{tabular}
    \caption{Dictionnaire -- Table projets}
    \label{tab:table_projets}
\end{table}

\begin{table}[H]
    \centering
    \small
    \begin{tabular}{lp{2cm}lp{6cm}}
        \toprule
        \rowcolor{primary!10}
        \textbf{Champ} & \textbf{Type} & \textbf{Contrainte} & \textbf{Description} \\
        \midrule
        id             & INT           & PK, AI              & Identifiant unique de la salle. \\
        nom            & VARCHAR       & NOT NULL            & Nom ou numéro de la salle. \\
        batiment       & VARCHAR       & -                   & Localisation géographique. \\
        capacite       & INT           & -                   & Nombre de places assises. \\
        \bottomrule
    \end{tabular}
    \caption{Dictionnaire -- Table salles}
    \label{tab:table_salles}
\end{table}

\paragraph{Soutenances et Évaluations.}

\begin{table}[H]
    \centering
    \small
    \begin{tabular}{lp{2cm}lp{6cm}}
        \toprule
        \rowcolor{primary!10}
        \textbf{Champ} & \textbf{Type} & \textbf{Contrainte} & \textbf{Description} \\
        \midrule
        id             & INT           & PK, AI              & Identifiant unique de la soutenance. \\
        date\_passage  & DATE          & NOT NULL            & Date prévue pour la séance. \\
        heure\_debut   & TIME          & NOT NULL            & Heure précise du début. \\
        duree          & INT           & -                   & Durée prévue en minutes. \\
        statut         & ENUM          & NOT NULL            & État (PLANIFIEE, CONVOQUEE, EVALUEE...). \\
        salle\_id      & INT           & FK                  & Référence à la salle occupée. \\
        projet\_id     & INT           & FK                  & Référence au projet présenté. \\
        session\_id    & INT           & FK                  & Référence à la session globale. \\
        \bottomrule
    \end{tabular}
    \caption{Dictionnaire -- Table soutenances}
    \label{tab:table_soutenances}
\end{table}

\begin{table}[H]
    \centering
    \small
    \begin{tabular}{lp{2cm}lp{6cm}}
        \toprule
        \rowcolor{primary!10}
        \textbf{Champ} & \textbf{Type} & \textbf{Contrainte} & \textbf{Description} \\
        \midrule
        soutenance\_id & INT           & PK, FK              & Référence à la soutenance. \\
        enseignant\_id & INT           & PK, FK              & Référence à l'enseignant membre. \\
        role\_jury     & ENUM          & NOT NULL            & Rôle (PRESIDENT, RAPPORTEUR, EXAMINATEUR). \\
        present        & BOOLEAN       & -                   & Confirmation de présence. \\
        \bottomrule
    \end{tabular}
    \caption{Dictionnaire -- Table membres\_jury}
    \label{tab:table_membres_jury}
\end{table}

\begin{table}[H]
    \centering
    \small
    \begin{tabular}{lp{2cm}lp{6cm}}
        \toprule
        \rowcolor{primary!10}
        \textbf{Champ} & \textbf{Type} & \textbf{Contrainte} & \textbf{Description} \\
        \midrule
        id             & INT           & PK, AI              & Identifiant de la note individuelle. \\
        note           & FLOAT         & NOT NULL            & Note attribuée par un membre du jury. \\
        observations   & TEXT          & -                   & Remarques du membre évaluateur. \\
        validee        & BOOLEAN       & DEFAULT FALSE       & Note confirmée. \\
        soutenance\_id & INT           & FK                  & Référence à la soutenance évaluée. \\
        enseignant\_id & INT           & FK                  & Membre du jury ayant saisi la note. \\
        \bottomrule
    \end{tabular}
    \caption{Dictionnaire -- Table evaluations}
    \label{tab:table_evaluations}
\end{table}

\begin{table}[H]
    \centering
    \small
    \begin{tabular}{lp{2cm}lp{6cm}}
        \toprule
        \rowcolor{primary!10}
        \textbf{Champ} & \textbf{Type} & \textbf{Contrainte} & \textbf{Description} \\
        \midrule
        id             & INT           & PK, AI              & Identifiant du résultat consolidé. \\
        moyenne\_finale & FLOAT        & NOT NULL            & Moyenne des notes validées (RG6). \\
        decision       & ENUM          & NOT NULL            & Décision finale (ADMIS, AJOURNE...). \\
        observations\_finales & TEXT   & -                   & Remarques finales du jury. \\
        validee        & BOOLEAN       & DEFAULT FALSE       & Verrouillage du résultat. \\
        soutenance\_id & INT           & FK, UNIQUE          & Soutenance concernée. \\
        \bottomrule
    \end{tabular}
    \caption{Dictionnaire -- Table resultats\_soutenance}
    \label{tab:table_resultats}
\end{table}

La table \texttt{evaluations} conserve la traçabilité des notes par membre du jury. La table \texttt{resultats\_soutenance} porte le résultat unique affiché à l'étudiant.

\paragraph{Documents, disponibilités et sécurité opérationnelle.}

\begin{table}[H]
    \centering
    \small
    \begin{tabular}{lp{2cm}lp{6cm}}
        \toprule
        \rowcolor{primary!10}
        \textbf{Champ} & \textbf{Type} & \textbf{Contrainte} & \textbf{Description} \\
        \midrule
        id             & INT           & PK, AI              & Identifiant de la convocation. \\
        type           & ENUM          & NOT NULL            & Type (ETUDIANT ou JURY). \\
        chemin\_fichier & VARCHAR      & -                   & Emplacement du PDF généré. \\
        envoyee        & BOOLEAN       & DEFAULT FALSE       & Envoi effectué. \\
        date\_generation & DATETIME    & NOT NULL            & Date de génération. \\
        soutenance\_id & INT           & FK                  & Soutenance concernée. \\
        utilisateur\_id & INT          & FK                  & Destinataire. \\
        \bottomrule
    \end{tabular}
    \caption{Dictionnaire -- Table convocations}
    \label{tab:table_convocations}
\end{table}

\begin{table}[H]
    \centering
    \small
    \begin{tabular}{lp{2cm}lp{6cm}}
        \toprule
        \rowcolor{primary!10}
        \textbf{Champ} & \textbf{Type} & \textbf{Contrainte} & \textbf{Description} \\
        \midrule
        id             & INT           & PK, AI              & Identifiant de l'indisponibilité. \\
        date\_debut    & DATE          & NOT NULL            & Premier jour indisponible. \\
        date\_fin      & DATE          & NOT NULL            & Dernier jour indisponible. \\
        heure\_debut   & TIME          & -                   & Début du créneau si partiel. \\
        heure\_fin     & TIME          & -                   & Fin du créneau si partiel. \\
        motif          & VARCHAR       & -                   & Motif ou commentaire. \\
        enseignant\_id & INT           & FK                  & Enseignant concerné. \\
        \bottomrule
    \end{tabular}
    \caption{Dictionnaire -- Table indisponibilites}
    \label{tab:table_indisponibilites}
\end{table}

\begin{table}[H]
    \centering
    \small
    \begin{tabular}{lp{2cm}lp{6cm}}
        \toprule
        \rowcolor{primary!10}
        \textbf{Champ} & \textbf{Type} & \textbf{Contrainte} & \textbf{Description} \\
        \midrule
        id             & INT           & PK, AI              & Identifiant de la notification. \\
        sujet          & VARCHAR       & NOT NULL            & Objet du message. \\
        contenu        & TEXT          & NOT NULL            & Corps du message. \\
        statut         & ENUM          & NOT NULL            & EN\_ATTENTE, ENVOYEE ou ECHEC. \\
        date\_envoi    & DATETIME      & -                   & Date effective d'envoi. \\
        utilisateur\_id & INT          & FK                  & Destinataire. \\
        soutenance\_id & INT           & FK                  & Soutenance liée. \\
        \bottomrule
    \end{tabular}
    \caption{Dictionnaire -- Table notifications}
    \label{tab:table_notifications}
\end{table}

\begin{table}[H]
    \centering
    \small
    \begin{tabular}{lp{2cm}lp{6cm}}
        \toprule
        \rowcolor{primary!10}
        \textbf{Champ} & \textbf{Type} & \textbf{Contrainte} & \textbf{Description} \\
        \midrule
        id             & INT           & PK, AI              & Identifiant du jeton. \\
        token          & VARCHAR       & UNIQUE, NOT NULL    & Valeur aléatoire du jeton sécurisé. \\
        expire\_le     & DATETIME      & NOT NULL            & Date limite d'utilisation. \\
        utilise        & BOOLEAN       & DEFAULT FALSE       & Empêche la réutilisation. \\
        utilisateur\_id & INT          & FK                  & Compte utilisateur concerné. \\
        \bottomrule
    \end{tabular}
    \caption{Dictionnaire -- Table reset\_tokens}
    \label{tab:table_reset_tokens}
\end{table}

\subsection{Conception de l'interface utilisateur}
\label{ch3:conception_ihm}

\subsubsection{Plan de navigation}
\label{ch3:sitemap}

Le diagramme suivant présente le sitemap de l'application et les points d'entrée selon les rôles utilisateurs.

\begin{figure}[H]
    \centering
    \pumlsvg[\textwidth]{nav_sitemap}
    \caption{Sitemap et plan de navigation}
    \label{fig:nav_sitemap}
\end{figure}

\subsubsection{Prototypes d'interface}
\label{ch3:prototypes}

Cette section présente les maquettes graphiques des principales interfaces de l'application. Le design privilégie la sobriété, l'efficacité et une navigation intuitive.

\paragraph{Interface d'accès.} L'écran de connexion constitue le point d'entrée unique du système. Il assure la sécurité des accès et redirige chaque utilisateur vers son espace dédié selon son rôle.

\begin{figure}[H]
    \centering
    
\includegraphics[width=\textwidth]{../../assets/prototype/login.png}
    \caption{Interface de connexion}
    \label{fig:login}
\end{figure}

\paragraph{Tableau de bord Administrateur.} Cette interface offre une vision globale sur l'infrastructure. Elle permet de gérer le cycle de vie des comptes utilisateurs et de configurer les ressources logistiques.

\begin{figure}[H]
    \centering
    \includegraphics[width=0.9\textwidth]{../../assets/prototype/admin-dashboard.png}
    \caption{Interface : Tableau de bord Administrateur}
    \label{fig:admin_dashboard}
\end{figure}

\paragraph{Tableau de bord Coordinateur.} Centre névralgique de la planification, cette vue permet au coordinateur d'importer les données, d'orchestrer les sessions et de visualiser l'état d'avancement.

\begin{figure}[H]
    \centering
    \includegraphics[width=0.9\textwidth]{../../assets/prototype/coordinator-dashboard.png}
    \caption{Interface : Tableau de bord Coordinateur}
    \label{fig:coordinator_dashboard}
\end{figure}

\paragraph{Espace Enseignant.} Cette interface permet aux enseignants de consulter leur planning de jury, de télécharger leurs convocations et d'accéder aux formulaires de notation.

\begin{figure}[H]
    \centering
    
\includegraphics[width=0.9\textwidth]{../../assets/prototype/teacher-dashboard.png}
    \caption{Interface : Tableau de bord Enseignant}
    \label{fig:teacher_dashboard}
\end{figure}

\paragraph{Espace Étudiant.} L'étudiant dispose d'une interface épurée affichant les informations cruciales : date, heure, salle et composition du jury, avec téléchargement de la convocation.

\begin{figure}[H]
    \centering
    
\includegraphics[width=0.9\textwidth]{../../assets/prototype/student-dashboard.png}
    \caption{Interface : Tableau de bord Étudiant}
    \label{fig:student_dashboard}
\end{figure}

\subsubsection{Modèle de convocation officielle}
\label{ch3:convocation_modele}

La convocation est le document administratif pivot du système. Elle constitue l'acte officiel qui déclenche la phase d'exécution des soutenances. L'automatisation complète de ce document via le service dédié apporte plusieurs garanties :
\begin{itemize}
    \item \textbf{Intégrité des données :} les informations sont extraites directement de la base après validation des conflits, éliminant les erreurs humaines.
    \item \textbf{Valeur administrative :} justificatif officiel d'absence pour l'étudiant et convocation de mission pour les enseignants.
    \item \textbf{Réactivité :} génération instantanée dès le verrouillage de la session par le coordinateur.
\end{itemize}

\begin{figure}[H]
    \centering
    \includesvg[width=\textwidth, inkscapelatex=false]{../../assets/prototype/convocation.svg}
    \caption{Spécimen de convocation officielle (Format A4)}
    \label{fig:convocation}
\end{figure}

\subsection{Conclusion}
\label{ch3:conclusion}

Ce chapitre a présenté l'analyse complète des besoins du système, depuis la définition des acteurs et des fonctionnalités jusqu'à la modélisation UML détaillée. Les règles de gestion qui encadrent le système ont été formalisées, les sept modules fonctionnels ont été décrits, et l'architecture technique -- couches MVC, composants, sécurité JWT -- a été spécifiée. Le modèle de données relationnel a été défini avec son dictionnaire complet. Enfin, la conception de l'interface utilisateur a été illustrée à travers les maquettes des écrans principaux et le modèle de convocation.

% ═══════════════════════════════════════════════════════════════════════════════
%  CHAPITRE 4 : RÉALISATION ET DÉVELOPPEMENT
% ═══════════════════════════════════════════════════════════════════════════════
\newpage
\section{Réalisation et développement}
\label{ch4}

\subsection{Introduction}
\label{ch4:intro}

Ce chapitre présente l'implémentation technique du système. Il décrit l'environnement de développement, l'architecture détaillée du code, l'implémentation des différents modules fonctionnels ainsi que les difficultés rencontrées et les solutions apportées.

\subsection{Environnement de développement}
\label{ch4:environnement}

Le développement a été réalisé avec les outils et configurations suivants :
\begin{itemize}[label=\textcolor{primary}{\textbullet}]
    \item \textbf{Système d'exploitation :} Windows 11 / Linux.
    \item \textbf{IDE Frontend :} VS Code avec extensions React, Tailwind CSS, ESLint.
    \item \textbf{IDE Backend :} IntelliJ IDEA Ultimate avec plugins Spring Boot, JPA, Maven.
    \item \textbf{Gestion de version :} Git et GitHub pour le versionnement du code source.
    \item \textbf{Test API :} Postman pour la validation des endpoints REST.
\end{itemize}

\subsection{Implémentation du backend (Spring Boot)}
\label{ch4:backend}

Le backend est structuré selon l'architecture Spring Boot en couches :
\begin{itemize}[label=\textcolor{primary}{\textbullet}]
    \item \textbf{Controllers :} exposent les endpoints REST pour chaque module.
    \item \textbf{Services :} contiennent la logique métier et les règles de gestion.
    \item \textbf{Repositories :} couche d'accès aux données via Spring Data JPA.
    \item \textbf{Entities :} classes Java représentant les tables de la base de données.
    \item \textbf{Security :} configuration Spring Security avec filtres JWT.
\end{itemize}

\subsection{Implémentation du frontend (React)}
\label{ch4:frontend}

Le frontend est développé avec React et Tailwind CSS :
\begin{itemize}[label=\textcolor{primary}{\textbullet}]
    \item \textbf{Composants :} architecture modulaire avec des composants réutilisables.
    \item \textbf{Routage :} React Router pour la navigation entre les vues.
    \item \textbf{État :} React Context et hooks pour la gestion d'état.
    \item \textbf{API :} Axios pour les appels HTTP vers le backend.
    \item \textbf{UI :} Shadcn UI pour les composants d'interface standardisés.
\end{itemize}

\subsubsection{Architecture détaillée du backend}
\label{ch4:backend_arch}

Le backend Spring Boot adopte une architecture en paquets métier strictement séparée par rôle fonctionnel, garantissant une faible couplage et une haute cohésion.

\begin{figure}[H]
    \centering
    \pumlsvg[0.8\textwidth]{pkg_backend_structure}
    \caption{Architecture en paquets du backend Spring Boot}
    \label{fig:backend_packages}
\end{figure}

Le paquetage racine `com.system_gestion_soutenance.api` est organisé comme suit :

\begin{itemize}[label=\textcolor{primary}{\textbullet}]
    \item \texttt{auth/} : Gestion de l'authentification (Login, JWT, mot de passe oublié)
    \item \texttt{user/} : Hiérarchie des utilisateurs avec héritage JOINED (\textbf{User} $\rightarrow$ Student/Teacher/Coordinator)
    \item \texttt{admin/} : Fonctionnalités réservées à l'administrateur
    \begin{itemize}[label=\scriptsize$\circ$]
        \item \texttt{session/} : Sessions globales académiques
        \item \texttt{defensesession/} : Sessions de soutenance spécifiques
        \item \texttt{room/} : Gestion des salles
        \item \texttt{department/} : Départements pédagogiques
        \item \texttt{config/} : Configuration système (niveaux, filières, grades, templates jury, emails, documents)
        \item \texttt{audit/} : Journalisation des actions
        \item \texttt{stats/} : Statistiques globales
    \end{itemize}
    \item \texttt{coordinator/} : Espace coordinateur
    \begin{itemize}[label=\scriptsize$\circ$]
        \item \texttt{project/} : Gestion des projets et sujets
        \item \texttt{group/} : Groupes d'étudiants
        \item \texttt{jury/} : Composition des jurys
        \item \texttt{schedule/} : Planification et créneaux
        \item \texttt{conflict/} : \textbf{Détection de conflits} (8 vérifications)
        \item \texttt{unavailability/} : Gestion des indisponibilités enseignants
    \end{itemize}
    \item \texttt{teacher/} : Espace enseignant/membre de jury
    \begin{itemize}[label=\scriptsize$\circ$]
        \item \texttt{evaluation/} : Saisie des notes et commentaires
        \item \texttt{schedule/} : Consultation du planning personnel
    \end{itemize}
    \item \texttt{student/} : Espace étudiant
    \begin{itemize}[label=\scriptsize$\circ$]
        \item \texttt{defense/} : Informations de soutenance
        \item \texttt{document/} : Dépôt des rapports
        \item \texttt{convocation/} : Téléchargement convocation
    \end{itemize}
    \item \texttt{notification/} : Notifications internes et emails
    \item \texttt{common/} : Configurations transverses
    \begin{itemize}[label=\scriptsize$\circ$]
        \item \texttt{config/} : SecurityConfig, OpenApiConfig
    \end{itemize}
\end{itemize}

\subsubsection{Architecture détaillée du frontend}
\label{ch4:frontend_arch}

Le frontend React (TypeScript) est structuré en couches logiques bien définies :

\begin{figure}[H]
    \centering
    \begin{minipage}{\linewidth}
        \centering
        \textit{[Placeholder : Diagramme d'architecture frontend]}
        \vspace{3cm}
    \end{minipage}
    \caption{Architecture du frontend React}
    \label{fig:frontend_arch}
\end{figure}

Structure des répertoires principaux :

\begin{itemize}[label=\textcolor{primary}{\textbullet}]
    \item \texttt{pages/} : 25+ pages organisées par rôle (Admin, Coordinator, Teacher, Student, Print)
    \item \texttt{components/} : Composants réutilisables
    \begin{itemize}[label=\scriptsize$\circ$]
        \item \texttt{ui/} : 39 composants primitifs Shadcn/Radix (button, card, dialog, data-table, sidebar...)
        \item \texttt{academic/} : 5 composants métier (ProjectDialog, CreateJuryDialog, ConvocationDialog...)
        \item \texttt{print/} : 8 composants dédiés à l'impression (Convocation, EvaluationSheet, AttendanceList, ProcesVerbal...)
    \end{itemize}
    \item \texttt{hooks/} : Hooks React Query personnalisés par rôle
    \begin{itemize}[label=\scriptsize$\circ$]
        \item \texttt{use-admin-queries.ts}, \texttt{use-coordinator-queries.ts}, etc.
    \end{itemize}
    \item \texttt{lib/} : Logique métier et couche API
    \begin{itemize}[label=\scriptsize$\circ$]
        \item \texttt{api-core.ts} : Wrapper fetch centralisé avec Bearer token
        \item \texttt{conflict-engine.ts} : 8 vérifications de conflits (côté client)
        \item \texttt{validations.ts} : Schémas Zod pour tous les formulaires
    \end{itemize}
    \item \texttt{contexts/} : AuthContext (état d'authentification global)
    \item \texttt{mocks/} : MSW (Mock Service Worker) pour le développement offline
    \begin{itemize}[label=\scriptsize$\circ$]
        \item Base de données in-memory complète
        \item Handlers pour tous les endpoints
    \end{itemize}
\end{itemize}

\subsection{Implémentation des modules fonctionnels}
\label{ch4:modules_impl}

Cette section détaille l'implémentation technique des sept modules fonctionnels définis dans le chapitre de conception.

\subsubsection{Module 1 : Authentification et gestion des rôles}

L'authentification repose sur le standard \textbf{JWT (JSON Web Token)} avec une architecture stateless.

\begin{figure}[H]
    \centering
    \pumlsvg[0.9\textwidth]{seq_auth}
    \caption{Diagramme de séquence : Processus d'authentification JWT}
    \label{fig:auth_seq_impl}
\end{figure}

Flux d'authentification :
\begin{enumerate}[label=\textcolor{primary}{\arabic*}.]
    \item L'utilisateur soumet \texttt{email + password}
    \item Backend: \texttt{BCryptPasswordEncoder} valide le hash
    \item \texttt{JwtTokenProvider} génère un token signé (HMAC-SHA256) valide 2 heures
    \item Claims inclus : \texttt{sub=userId}, \texttt{role=ADMIN|COORDINATOR|TEACHER|STUDENT}
    \item Le frontend stocke le token dans \texttt{localStorage}
    \item Chaque requête ultérieure inclut l'en-tête \texttt{Authorization: Bearer <token>}
\end{enumerate}

Sécurité côté backend (\texttt{SecurityConfig.java}) :
\begin{itemize}[label=\textcolor{primary}{\textbullet}]
    \item \textbf{Stateless} : Aucune session côté serveur
    \item \textbf{Role-based access control} :
    \begin{itemize}[label=\scriptsize$\circ$]
        \item \texttt{/api/admin/**} $\rightarrow$ \texttt{ROLE\_ADMIN}
        \item \texttt{/api/coordinator/**} $\rightarrow$ \texttt{ROLE\_COORDINATOR}
        \item \texttt{/api/teacher/**} $\rightarrow$ \texttt{ROLE\_TEACHER}
        \item \texttt{/api/student/**} $\rightarrow$ \texttt{ROLE\_STUDENT}
    \end{itemize}
    \item \textbf{Endpoints publics} : \texttt{/api/login}, \texttt{/api/auth/**}, \texttt{/swagger-ui/**}, \texttt{/h2-console/**}
    \item Custom handlers pour \texttt{401 Unauthorized} et \texttt{403 Access Denied}
\end{itemize}

\subsubsection{Module 2 : Gestion académique}

Ce module centralise la gestion des étudiants, enseignants, projets, groupes et configurations académiques.

\begin{itemize}[label=\textcolor{primary}{\textbullet}]
    \item \textbf{Hiérarchie User} : Héritage \texttt{JOINED} avec discriminant par rôle
    \item \textbf{Entités de configuration} : \texttt{Major} (filière), \texttt{Level} (niveau), \texttt{Grade}, \texttt{Department}
    \item \textbf{Projets} : Association \texttt{ManyToMany} avec \texttt{Student} via \texttt{Group}
    \item \textbf{Import CSV} : Prévu pour utilisateurs et salles (bulk create)
\end{itemize}

\subsubsection{Module 3 : Planification}

Le cœur de l'application : création des créneaux, réservation des salles, génération automatique optionnelle.

\begin{figure}[H]
    \centering
    \pumlsvg[\textwidth]{seq_planning}
    \caption{Diagramme de séquence : Processus de planification}
    \label{fig:planning_seq_impl}
\end{figure}

Entité clé : \texttt{SlotAssignment}
\begin{itemize}[label=\textcolor{primary}{\textbullet}]
    \item \texttt{date}, \texttt{time} : Créneau horaire
    \item \texttt{projectId} : Projet associé
    \item \texttt{room} : Salle réservée
\end{itemize}

Fonctionnalités implémentées :
\begin{itemize}[label=\textcolor{primary}{\textbullet}]
    \item \texttt{GET/POST /api/coordinator/schedule} : Consultation et modification
    \item \texttt{POST /auto-generate} : Algorithme itératif qui peuple le planning
    \item \texttt{POST /publish} : Verrouille le planning et crée des notifications
\end{itemize}

\subsubsection{Module 4 : Gestion des jurys}

Composition des jurys avec respect strict des Règles de Gestion RG4-RG5.

\begin{figure}[H]
    \centering
    \pumlsvg[\textwidth]{seq_jury_assignment}
    \caption{Diagramme de séquence : Affectation d'un jury}
    \label{fig:jury_seq_impl}
\end{figure}

Modèle de données :
\begin{itemize}[label=\textcolor{primary}{\textbullet}]
    \item \texttt{Jury} : Entité container
    \item \texttt{JuryMember} : Ligne d'association
    \begin{itemize}[label=\scriptsize$\circ$]
        \item \texttt{roleName} : \texttt{"PRESIDENT"}, \texttt{"RAPPORTEUR"}, \texttt{"EXAMINATEUR"}, \texttt{"SUPERVISOR"}
        \item \texttt{teacher} : ManyToOne vers Teacher
    \end{itemize}
    \item \texttt{JuryRoleTemplate} : Templates configurables par type de soutenance
\end{itemize}

\paragraph{Moteur de détection de conflits (8 vérifications)}

Le service \texttt{ConflictDetectionService} effectue systématiquement ces vérifications :

\begin{table}[H]
    \centering
    \small
    \begin{tabular}{lp{10cm}}
        \toprule
        \rowcolor{primary!10}
        \textbf{Conflit} & \textbf{Vérification} \\
        \midrule
        \texttt{project\_already\_scheduled} & Un projet ne peut avoir deux créneaux \\
        \texttt{slot\_occupied} & Même date+heure+salle déjà prise \\
        \texttt{room\_capacity} & Capacité de la salle respectée \\
        \texttt{date\_out\_of\_bounds} & Date dans la session globale \\
        \texttt{teacher\_double\_booked} & Enseignant pas deux soutenances simultanées \\
        \texttt{supervisor\_conflict} & Encadrant disponible \\
        \texttt{break\_violation} & Pause entre deux soutenances \\
        \texttt{teacher\_unavailable} & Respect des indisponibilités déclarées \\
        \bottomrule
    \end{tabular}
    \caption{Types de conflits détectés}
    \label{tab:conflict_types}
\end{table}

Ce moteur est implémenté \textbf{à la fois} backend (\texttt{ConflictDetectionService.java}) et frontend (\texttt{conflict-engine.ts}) pour une validation en temps réel lors du drag-and-drop.

\subsubsection{Module 5 : Évaluation}

Saisie des notes par les membres du jury.

\begin{figure}[H]
    \centering
    \pumlsvg[\textwidth]{seq_evaluation}
    \caption{Diagramme de séquence : Saisie d'une évaluation}
    \label{fig:eval_seq_impl}
\end{figure}

Entité \texttt{Evaluation} :
\begin{itemize}[label=\textcolor{primary}{\textbullet}]
    \item \texttt{teacherId} : Auteur de l'évaluation
    \item \texttt{projectId} : Projet concerné
    \item \texttt{role} : Rôle dans le jury (Président, Rapporteur...)
    \item \texttt{score} : Note numérique
    \item \texttt{comment} : Observations textuelles
    \item \texttt{status} : \texttt{DRAFT} / \texttt{SUBMITTED} (soumission unique)
\end{itemize}

Les coefficients de pondération par rôle sont configurables dans \texttt{DefenseSession.evaluationCoefficients}.

\subsubsection{Module 6 : Génération documentaire}

Deux approches complémentaires :

\paragraph{Backend : Convocation endpoint}
Endpoint \texttt{GET /api/student/convocation} :
\begin{itemize}[label=\textcolor{primary}{\textbullet}]
    \item Actuellement retourne un placeholder texte
    \item Prêt pour implémentation PDF (iText ou Apache PDFBox)
\end{itemize}

\paragraph{Frontend : Composants Print}

7 pages React dédiées sous \texttt{/print/*} :
\begin{itemize}[label=\textcolor{primary}{\textbullet}]
    \item \texttt{/print/evaluation-sheet}
    \item \texttt{/print/attendance-list}
    \item \texttt{/print/jury-convocation}
    \item \texttt{/print/schedule}
    \item \texttt{/print/proces-verbal}
    \item \texttt{/print/student-convocation}
    \item \texttt{/print/certificate}
\end{itemize}

Ces composants utilisent des styles CSS \texttt{@media print} pour une sortie PDF de qualité via le navigateur (Ctrl+P).

\subsubsection{Module 7 : Tableau de bord}

Tableaux de bord différenciés par rôle :

\begin{itemize}[label=\textcolor{primary}{\textbullet}]
    \item \textbf{Admin Dashboard} : Statistiques globales, compteurs utilisateurs, graphiques Recharts
    \item \textbf{Coordinator Dashboard} : Vue d'ensemble session, avancement planification
    \item \textbf{Teacher Dashboard} : Prochaines soutenances, évaluations en attente
    \item \textbf{Student Dashboard} : Informations de sa soutenance, composition jury
\end{itemize}

\subsection{Extraits de code représentatifs}
\label{ch4:code_snippets}

\subsubsection{1. Génération JWT (Backend)}

\begin{minted}[frame=single, linenos, breaklines, fontsize=\small]{java}
// JwtTokenProvider.java
@Service
public class JwtTokenProvider {

    @Value("${app.jwt-secret:s3cr3t-k3y-f0r-d3f3ns3-m4n4g3m3nt-syst3m-2026}")
    private String jwtSecret;

    @Value("${app.jwt-expiration-ms:7200000}") // 2 hours
    private int jwtExpirationMs;

    public String generateToken(Integer userId, Role role) {
        Date now = new Date();
        Date expiryDate = new Date(now.getTime() + jwtExpirationMs);

        return Jwts.builder()
                .subject(userId.toString())
                .claim("role", role.name())
                .issuedAt(now)
                .expiration(expiryDate)
                .signWith(getSignInKey())
                .compact();
    }

    private SecretKey getSignInKey() {
        byte[] keyBytes = Decoders.BASE64.decode(
            java.util.Base64.getEncoder().encodeToString(jwtSecret.getBytes())
        );
        return Keys.hmacShaKeyFor(keyBytes);
    }
}
\end{minted}

\subsubsection{2. Vérification de conflit : Double réservation enseignant}

\begin{minted}[frame=single, linenos, breaklines, fontsize=\small]{java}
// ConflictDetectionService.java - extrait
@Service
public class ConflictDetectionService {

    public List<ConflictResult> validateSlotAssignment(SlotAssignment slot) {
        List<ConflictResult> conflicts = new ArrayList<>();

        // Récupère le projet et son jury
        Optional<Project> project = projectRepository.findById(slot.getProjectId());
        if (project.isEmpty()) return conflicts;

        // Vérifie chaque membre du jury
        List<JuryMember> members = juryMemberRepository.findByJury_ProjectId(slot.getProjectId());
        
        for (JuryMember member : members) {
            LocalTime slotStart = slot.getTime();
            LocalTime slotEnd = slotStart.plusMinutes(DEFAULT_DURATION);
            
            // Cherche d'autres créneaux pour ce même enseignant à la même date
            List<SlotAssignment> teacherSlots = slotAssignmentRepository
                .findByDateAndJuryMembersContaining(slot.getDate(), member.getTeacher());
            
            for (SlotAssignment existing : teacherSlots) {
                if (existing.getId().equals(slot.getId())) continue;
                
                LocalTime existingStart = existing.getTime();
                LocalTime existingEnd = existingStart.plusMinutes(DEFAULT_DURATION);
                
                // Vérifie chevauchement
                if (timesOverlap(slotStart, slotEnd, existingStart, existingEnd)) {
                    conflicts.add(ConflictResult.builder()
                        .type("teacher_double_booked")
                        .severity(Severity.ERROR)
                        .teacher(member.getTeacher())
                        .conflictingSlot(existing)
                        .message("Enseignant " + member.getTeacher().getUser().getFullName() + 
                                 " déjà affecté à " + existing.getTitle() + 
                                 " (" + existing.getTime() + ")")
                        .build());
                }
            }
        }
        return conflicts;
    }
}
\end{minted}

\subsubsection{3. React Query : Hook personnalisé (Frontend)}

\begin{minted}[frame=single, linenos, breaklines, fontsize=\small]{typescript}
// use-coordinator-queries.ts - extrait
import { useQuery, useMutation, useQueryClient } from '@tanstack/react-query';
import { 
  getProjects, createProject, updateProject, deleteProject,
  getSchedule, saveSchedule, validateSchedule, generateSchedule,
  type CreateProjectRequest, type SlotAssignment 
} from '@/lib/api-coordinator';

export function useProjects() {
  return useQuery({
    queryKey: ['coordinator', 'projects'],
    queryFn: getProjects,
    staleTime: 30_000 // 30s
  });
}

export function useCreateProject() {
  const queryClient = useQueryClient();
  
  return useMutation({
    mutationFn: (data: CreateProjectRequest) => createProject(data),
    onSuccess: () => {
      queryClient.invalidateQueries({ queryKey: ['coordinator', 'projects'] });
      queryClient.invalidateQueries({ queryKey: ['coordinator', 'stats'] });
    }
  });
}

export function useSchedule() {
  return useQuery({
    queryKey: ['coordinator', 'schedule'],
    queryFn: getSchedule
  });
}

export function useGenerateSchedule() {
  const queryClient = useQueryClient();
  
  return useMutation({
    mutationFn: () => generateSchedule(),
    onSuccess: () => {
      queryClient.invalidateQueries({ queryKey: ['coordinator', 'schedule'] });
    }
  });
}
\end{minted}

\subsection{Interfaces utilisateur}
\label{ch4:interfaces}

Cette section présente les principales vues de l'application.

\begin{table}[H]
    \centering
    \small
    \begin{tabular}{lp{6cm}p{5cm}}
        \toprule
        \rowcolor{primary!10}
        \textbf{Rôle} & \textbf{Pages principales} & \textbf{URL} \\
        \midrule
        \multirow{6}{*}{\textbf{Admin}}
            & Tableau de bord & \texttt{/admin} \\
            & Gestion utilisateurs & \texttt{/admin/users/*} \\
            & Salles & \texttt{/admin/rooms} \\
            & Sessions académiques & \texttt{/admin/sessions} \\
            & Configuration & \texttt{/admin/config} \\
            & Journal d'audit & \texttt{/admin/audit-logs} \\
        \midrule
        \multirow{6}{*}{\textbf{Coordinator}}
            & Tableau de bord & \texttt{/coordinator} \\
            & Projets \& Groupes & \texttt{/coordinator/projects} \\
            & \textbf{Defense Designer} & \texttt{/coordinator/schedule} \\
            & Jurys & \texttt{/coordinator/juries} \\
            & Conflits & \texttt{/coordinator/conflicts} \\
            & Notes & \texttt{/coordinator/grades} \\
        \midrule
        \multirow{4}{*}{\textbf{Teacher}}
            & Tableau de bord & \texttt{/teacher} \\
            & Planning & \texttt{/teacher/schedule} \\
            & Évaluations & \texttt{/teacher/evaluations} \\
            & Indisponibilités & \texttt{/teacher/unavailability} \\
        \midrule
        \multirow{3}{*}{\textbf{Student}}
            & Tableau de bord & \texttt{/student} \\
            & Groupe & \texttt{/student/group} \\
            & Documents & \texttt{/student/documents} \\
        \bottomrule
    \end{tabular}
    \caption{Architecture des pages par rôle}
    \label{tab:pages_par_role}
\end{table}

\subsubsection{Écran de connexion}

Point d'entrée unique pour tous les utilisateurs avec validation des identifiants et redirection selon le rôle.

\begin{figure}[H]
    \centering
    \begin{minipage}{0.7\textwidth}
        \centering
        \framebox[\linewidth][c]{\rule{0pt}{5cm} \textbf{[Placeholder : Capture d'écran Login]}}
    \end{minipage}
    \caption{Interface : Page de connexion}
    \label{fig:ui_login}
\end{figure}

\subsubsection{Defense Designer (Coordinateur)}

Vue centrale de planification avec système de glisser-déposer (DnD Kit) et validation de conflits en temps réel.

\begin{figure}[H]
    \centering
    \begin{minipage}{\linewidth}
        \centering
        \framebox[\linewidth][c]{\rule{0pt}{6cm} \textbf{[Placeholder : Capture d'écran Defense Designer]}}
    \end{minipage}
    \caption{Interface : Defense Designer (Planification)}
    \label{fig:ui_designer}
\end{figure}

Fonctionnalités clés de cette vue :
\begin{itemize}[label=\textcolor{primary}{\textbullet}]
    \item Calendrier par jour / par salle
    \item Projets non planifiés dans la barre latérale
    \item Drag \& drop pour assigner un projet à un créneau
    \item Indicateurs visuels de conflits (rouge = bloquant, orange = avertissement)
    \item Boutons : \textbf{Auto-generate}, \textbf{Validate}, \textbf{Publish}
\end{itemize}

\subsubsection{Tableau de bord Administrateur}

Vue globale avec compteurs et graphiques Recharts.

\begin{figure}[H]
    \centering
    \begin{minipage}{0.9\textwidth}
        \centering
        \framebox[\linewidth][c]{\rule{0pt}{5cm} \textbf{[Placeholder : Capture d'écran Admin Dashboard]}}
    \end{minipage}
    \caption{Interface : Tableau de bord Admin}
    \label{fig:ui_admin}
\end{figure}

\subsubsection{Formulaire d'évaluation (Enseignant)}

Interface de notation par membre de jury avec champ commentaire et coefficients.

\begin{figure}[H]
    \centering
    \begin{minipage}{0.7\textwidth}
        \centering
        \framebox[\linewidth][c]{\rule{0pt}{6cm} \textbf{[Placeholder : Capture d'écran Evaluation Form]}}
    \end{minipage}
    \caption{Interface : Formulaire d'évaluation}
    \label{fig:ui_evaluation}
\end{figure}

\subsubsection{Espace Étudiant}

Vue simplifiée affichant les informations de soutenance : date, heure, salle, composition du jury.

\begin{figure}[H]
    \centering
    \begin{minipage}{0.9\textwidth}
        \centering
        \framebox[\linewidth][c]{\rule{0pt}{5cm} \textbf{[Placeholder : Capture d'écran Student Dashboard]}}
    \end{minipage}
    \caption{Interface : Espace étudiant}
    \label{fig:ui_student}
\end{figure}

\subsubsection{Impression : Convocation}

Composant \texttt{/print/student-convocation} optimisé pour export PDF (Ctrl+P).

\begin{figure}[H]
    \centering
    \begin{minipage}{0.7\textwidth}
        \centering
        \framebox[\linewidth][c]{\rule{0pt}{8cm} \textbf{[Placeholder : Capture d'écran Convocation PDF]}}
    \end{minipage}
    \caption{Interface : Modèle de convocation (impression)}
    \label{fig:ui_convocation}
\end{figure}

Autres vues d'impression disponibles :
\begin{itemize}[label=\textcolor{primary}{\textbullet}]
    \item Feuille d'évaluation
    \item Liste de présence
    \item Planning global
    \item Procès-verbal
    \item Convocation jury
\end{itemize}

\subsection{Difficultés rencontrées et solutions}
\label{ch4:difficultes}

Les principales difficultés rencontrées lors du développement incluent la gestion optimisée des conflits horaires et l'équilibrage automatique de la charge des jurys, résolues par des algorithmes de vérification systématique et des heuristiques de répartition.

\subsection{Conclusion}
\label{ch4:conclusion}

Ce chapitre a présenté l'implémentation du système selon l'architecture définie au chapitre précédent. Le backend Spring Boot et le frontend React ont été structurés pour assurer la maintenabilité et l'évolutivité de l'application.

% ═══════════════════════════════════════════════════════════════════════════════
%  CHAPITRE 5 : TESTS ET VALIDATION
% ═══════════════════════════════════════════════════════════════════════════════
\newpage
\section{Tests et validation}
\label{ch5}

\subsection{Introduction}
\label{ch5:intro}

Ce chapitre présente la stratégie de test adoptée pour garantir la qualité et la fiabilité du système développé. Il détaille les tests fonctionnels et d'intégration, les scénarios de validation, ainsi que les critères de réussite du projet.

\subsection{Stratégie de test}
\label{ch5:strategie}

La stratégie de test adoptée prévoit trois niveaux complémentaires, avec une infrastructure technique déjà configurée :

\begin{itemize}[label=\textcolor{primary}{\textbullet}]
    \item \textbf{Tests unitaires et d'intégration (Backend)} : Infrastructure \textbf{JUnit 5} et \textbf{Spring Boot Test} configurée dans le projet Maven.
    \item \textbf{Tests d'intégration API} : Prévus via \textbf{Postman} avec collections d'endpoints.
    \item \textbf{Tests End-to-End (Frontend)} : Projet séparé \texttt{system-gestion-soutenance-e2e} avec \textbf{Playwright} configuré.
    \item \textbf{Mocks développement} : \textbf{MSW (Mock Service Worker)} implémenté côté frontend pour prototypage offline.
\end{itemize}

\begin{table}[H]
    \centering
    \small
    \begin{tabular}{lll}
        \toprule
        \rowcolor{primary!10}
        \textbf{Niveau de test} & \textbf{Framework} & \textbf{État} \\
        \midrule
        Tests unitaires Backend & JUnit 5 + Spring Boot Test & Infrastructure OK \\
        Tests API & Postman & Collections à créer \\
        Tests E2E & Playwright & Projet prêt \\
        Tests frontend offline & MSW & \textbf{Implémenté} (développement) \\
        \bottomrule
    \end{tabular}
    \caption{Infrastructure de test}
    \label{tab:test_infra}
\end{table}

\subsection{Scénarios de test}
\label{ch5:scenarios}

Les scénarios de test planifiés couvrent les fonctionnalités critiques du système. \textbf{À la date de rédaction de ce rapport, les tests automatisés ne sont pas encore implémentés} (seul un test placeholder \texttt{contextLoads()} existe dans le backend).

\begin{projectbox}{Tests manuels effectués}
    Toutes les fonctionnalités principales ont été validées \textbf{manuellement} pendant le développement :
    \begin{itemize}[label=\textcolor{primary}{\textbullet}, leftmargin=1cm]
        \item Connexion/déconnexion avec les 4 rôles
        \item Gestion CRUD complète (utilisateurs, salles, projets, jurys, créneaux)
        \item Planification avec détection de conflits en temps réel
        \item Saisie et consultation des évaluations
        \item Composants d'impression PDF via navigateur
    \end{itemize}
\end{projectbox}

Scénarios prévus pour l'automatisation future :

\begin{itemize}[label=\textcolor{primary}{\textbullet}]
    \item \textbf{Authentification :} connexion avec identifiants valides/invalides, restriction d'accès par rôle, expiration de session.
    \item \textbf{Gestion des utilisateurs :} création, modification et suppression de comptes pour chaque rôle.
    \item \textbf{Planification :} création d'une session globale, planification d'une soutenance individuelle, modification des créneaux.
    \item \textbf{Affectation des jurys :} composition d'un jury avec respect des règles RG4 et RG5, équilibrage de la charge.
    \item \textbf{Détection de conflits :} double affectation d'un enseignant, réservation simultanée d'une salle, double planification d'un étudiant.
    \item \textbf{Génération de documents :} production de convocations PDF, fiches d'évaluation et planning général.
    \item \textbf{Saisie des évaluations :} enregistrement des notes, calcul de la moyenne (RG6), validation du résultat final.
\end{itemize}

\subsection{Validation des fonctionnalités}
\label{ch5:validation}

La validation manuelle a été effectuée selon les critères suivants :

\begin{itemize}[label=\textcolor{primary}{\textbullet}]
    \item Les utilisateurs peuvent être gérés de manière sécurisée selon leur rôle.
    \item Les étudiants, enseignants et salles peuvent être enregistrés et modifiés sans erreur.
    \item Une soutenance peut être planifiée en respectant les contraintes temporelles et les règles de gestion.
    \item Un jury peut être affecté avec détection automatique des conflits.
    \item Les composants d'impression (convocations, évaluations) sont fonctionnels et exportables en PDF via le navigateur.
    \item Les évaluations et notes peuvent être saisies et enregistrées de manière persistante.
    \item Le planning global est consultable, filtrable et exploitable par tous les acteurs.
\end{itemize}

\subsection{Évaluation globale du système}
\label{ch5:evaluation}

L'évaluation globale du système repose actuellement sur les tests manuels effectués par les développeurs tout au long du cycle de développement.

Le \textbf{groupe pilote} -- coordinateur, enseignants et étudiants représentant les rôles décrits en Section \ref{ch3:acteurs} -- n'a pas encore été mobilisé. Cette étape de validation utilisateur est prévue pour les prochaines itérations du projet, afin de recueillir des retours sur l'ergonomie, la pertinence des fonctionnalités et la performance globale en conditions réelles.

\subsection{Limites du projet}
\label{ch5:limites}

Le projet présente certaines limites qui pourront être adressées dans des versions futures :

\begin{itemize}[label=\textcolor{primary}{\textbullet}]
    \item La planification optimisée des créneaux n'est pas entièrement automatisée ; elle repose sur les choix du coordinateur assisté par la détection de conflits.
    \item L'envoi des notifications par e-mail n'est pas implémenté dans la version initiale (mock seulement).
    \item L'import massif de données depuis Excel est limité à un format prédéfini.
    \item \textbf{Tests automatisés :} L'infrastructure est configurée (JUnit 5, Spring Boot Test, Playwright) mais les tests ne sont pas implémentés (seulement un test placeholder).
    \item \textbf{Génération PDF backend :} L'endpoint \texttt{/api/student/convocation} retourne actuellement un placeholder texte ; la génération PDF native via iText/PDFBox est à implémenter.
\end{itemize}

\subsection{Conclusion}
\label{ch5:conclusion}

La stratégie de test définie dans ce chapitre pose les bases pour garantir la conformité du système aux exigences du cahier des charges. L'infrastructure technique -- \textbf{JUnit 5, Spring Boot Test, MSW, Playwright} -- est en place et prête à accueillir les tests automatisés dans les prochaines itérations.

En attendant, une validation manuelle approfondie a été effectuée sur l'ensemble des fonctionnalités principales : authentification multi-rôles, planification avec détection de conflits, gestion des jurys, saisie des évaluations et composants d'impression.

% ═══════════════════════════════════════════════════════════════════════════════
%  CONCLUSION GÉNÉRALE ET PERSPECTIVES
% ═══════════════════════════════════════════════════════════════════════════════
\newpage
\section*{Conclusion Générale et Perspectives}
\addcontentsline{toc}{section}{Conclusion Générale et Perspectives}

\subsection*{Bilan du travail réalisé}
\addcontentsline{toc}{subsection}{Bilan du travail réalisé}

Ce projet de fin d'études nous a permis de concevoir et de développer une application web complète pour la gestion des soutenances et des jurys au sein de la Faculté des Sciences Ben M'Sick. Partant d'un constat partagé -- l'inefficacité des méthodes manuelles actuelles -- nous avons mené une analyse approfondie de l'existant, étudié les solutions disponibles sur le marché et défini un cahier des charges précis.

La conception du système a été réalisée en suivant une démarche rigoureuse de modélisation UML, couvrant les aspects fonctionnels, statiques et dynamiques. L'architecture technique retenue -- Spring Boot pour le backend, React pour le frontend, MySQL pour la persistance -- offre un socle robuste, évolutif et sécurisé.

Le système couvre l'ensemble du cycle de vie d'une soutenance, de la planification des créneaux et de l'affectation des jurys jusqu'à la saisie des évaluations et la génération automatique des documents administratifs, en passant par la détection des conflits.

\subsection*{Apports personnels et compétences acquises}
\addcontentsline{toc}{subsection}{Apports personnels et compétences acquises}

La réalisation de ce projet nous a permis de mettre en pratique les connaissances acquises durant notre formation en licence MIP, notamment en développement web, en bases de données et en génie logiciel. Nous avons approfondi notre maîtrise de technologies modernes telles que Spring Boot, React, JWT et MySQL.

Sur le plan méthodologique, ce projet nous a initiés à la gestion d'un projet informatique complet, depuis l'analyse des besoins jusqu'à la conception, en passant par la modélisation UML et la gestion de version avec Git.

\subsection*{Perspectives d'évolution}
\addcontentsline{toc}{subsection}{Perspectives d'évolution}

Plusieurs évolutions sont envisagées pour enrichir le système à moyen terme :
\begin{itemize}[label=\textcolor{primary}{\textbullet}]
    \item \textbf{Notifications enrichies :} modèles d'e-mails avancés, rappels automatiques et suivi détaillé des envois.
    \item \textbf{Import avancé :} chargement massif de données académiques depuis des fichiers Excel.
    \item \textbf{Planification optimisée :} génération automatique des plannings fondée sur des algorithmes de satisfaction de contraintes.
    \item \textbf{Visioconférence intégrée :} module dédié aux soutenances à distance.
    \item \textbf{Signature électronique :} des procès-verbaux et archivage numérique des soutenances passées.
    \item \textbf{Statistiques annuelles :} par filière, avec tableaux de bord enrichis.
    \item \textbf{Interopérabilité :} intégration avec le système d'information global de l'établissement.
\end{itemize}

% ═══════════════════════════════════════════════════════════════════════════════
%  BIBLIOGRAPHIE
% ═══════════════════════════════════════════════════════════════════════════════
\newpage
\section*{Bibliographie}
\addcontentsline{toc}{section}{Bibliographie}

\begin{thebibliography}{9}

\bibitem{spring}
Spring Boot Reference Documentation,
\url{https://docs.spring.io/spring-boot/docs/current/reference/html/}

\bibitem{react}
React Documentation,
\url{https://react.dev/}

\bibitem{jwt}
JSON Web Token (JWT) RFC 7519,
\url{https://datatracker.ietf.org/doc/html/rfc7519}

\bibitem{uml}
OMG Unified Modeling Language (UML) Specification,
\url{https://www.omg.org/spec/UML/}

\bibitem{mysql}
MySQL Documentation,
\url{https://dev.mysql.com/doc/}

\bibitem{shadcn}
Shadcn UI Documentation,
\url{https://ui.shadcn.com/}

\bibitem{tailwind}
Tailwind CSS Documentation,
\url{https://tailwindcss.com/docs}

\bibitem{git}
Pro Git Book,
\url{https://git-scm.com/book/}

\end{thebibliography}

\end{document}
